% !TEX root = ../main.tex
%
% Chapter 4. View II: the Borcherds--Kac--Moody denominator algebra
% \autcor on \La^{2,1}_{II}. Cartan / Chevalley / real and imaginary
% roots / Weyl vector / denominator identity / Borcherds-weight
% characteristic. Imported from Gritsenko--Nikulin (1997, 1998),
% Borcherds (1988, 1995, 1998), Gritsenko--Nikulin (1997, 1998), Kac (1990); cross-cited in Vol III.
%
% This chapter is the chiral shadow of the Beilinson cut for
% View II. It is target-side. No compact Hall claim, Pfaffian
% claim, or BPS Hilbert-space claim is made.

\chapter{The Borcherds--Kac--Moody denominator algebra
\texorpdfstring{$\autcor$}{g\_\{Delta\_5\}} on
\texorpdfstring{$\La^{2,1}_{II}$}{Lambda\^\{2,1\}\_II}}
\label{ch:view2-bkm}

The Igusa cusp form $\Delta_5$ admits a Borcherds--Kac--Moody
description as a denominator. The denominator algebra is
the generalized Kac--Moody Lie superalgebra $\autcor$ on the
Lorentzian root lattice $\La^{2,1}_{II}=\HypPlan\oplus[2]$ in its
type-II sublattice presentation, with three even real simple roots
$\delta_1,\delta_2,\delta_3$, an infinite imaginary simple-root datum
read from the K3 half-index $\phi_{0,1}$, and Weyl group
$W^{(2)}(\La^{2,1}_{II})$ generated by the three real reflections.
Its Weyl--Kac--Borcherds denominator is the level-five Borcherds--Kac--Moody
denominator
\[
 \operatorname{den}(\autcor)\;=\;64^{-1}\Delta_5(2Z),
\]
the doubled Igusa form normalized by the theta-constant constant
$[q^{1/2}r^{1/2}s^{1/2}]\Delta_5=64$. The Borcherds-weight
characteristic of the denominator is
$\kappa_{\mathrm{BKM}}(\Delta_5)=c_1(0)/2=10/2=5$, the weight of the
Gritsenko paramodular form $\Delta_5=\Phi_1$ at level $1$ —
equivalently $\Phi_{10}^{1/2}$ in Igusa's normalization, where
$\Phi_{10}=\chi_{10}=\Delta_5^2$.  The handle $\Phi_1$ in this chapter
denotes Gritsenko's paramodular $\Phi_t$ at $t=1$, distinct from
Igusa's weight-ten $\Phi_{10}$.

The chapter is target-side. The denominator algebra is the chiral
shadow of an open compact realization. No compact Hall bracket,
protected Pfaffian, or BPS Hilbert space is claimed; those are the
conditional contribution of Chapter~\ref{chap:dirac-igusa-pro-object}
under the four obstructions $(D0)$, $(O1)$, $(O1)^+$, $(O2)$.

The three named instances of $\Delta_5$ are tagged distinctly. In
this chapter the canonical handle is the Borcherds--Kac--Moody
denominator,
\[
 \operatorname{den}(\autcor)\;=\;64^{-1}\Delta_5(2Z).
\]
The automorphic section $\mathcal D_X=\Delta_5\in M_5(\Sp_4(\Z),
\nu_{\Delta_5})$ of Chapter~\ref{ch:view1-automorphic} is View~I; the compact
protected Pfaffian $\operatorname{Pf}_{\mathrm{prot}}(\mathfrak D_X^{DI})$ of
Chapter~\ref{chap:dirac-igusa-pro-object} is View~III, conditional. The bare
handle $\Delta_5$ is the working symbol for the Siegel cusp form of
weight five and never substitutes for any of the three.

\section{The hyperbolic lattice
\texorpdfstring{$\La^{2,1}=\HypPlan\oplus[2]$}{Lambda\^\{2,1\} = U
oplus [2]} and its type-II sublattice}
\label{sec:bkm-lattice}

Fix the rank-three Lorentzian lattice
\begin{equation}
 \label{eq:bkm-lambda21-def}
 \La^{2,1}\;=\;\HypPlan\oplus[2]\;=\;\Z f_2\oplus\Z f_3\oplus\Z f_{-2},
\end{equation}
where $\HypPlan$ is the hyperbolic plane and $[2]$ is the rank-one
lattice generated by a single vector of square $+2$. The bilinear form
on the chosen basis is
\begin{equation}
 \label{eq:bkm-lambda21-form}
 (f_2,f_{-2})\;=\;-1,\qquad
 (f_2,f_2)\;=\;(f_{-2},f_{-2})\;=\;0,\qquad
 (f_3,f_3)\;=\;2,
\end{equation}
with all other pairings between basis vectors zero. The lattice is
even, signature $(2,1)$, and primitively embedded in the rank-five
Lorentzian lattice $\La^{3,2}\simeq\HypPlan\oplus\HypPlan\oplus[2]$ of
Section~\ref{ssec:exceptional-iso}.

\begin{definition}[Type-II sublattice]
\label{def:bkm-type-ii-sublattice}
The type-II sublattice of $\La^{2,1}$ is
\begin{equation}
 \label{eq:bkm-type-ii-def}
 \La^{2,1}_{II}
 \;=\;
 \{mf_2+lf_3+nf_{-2}\in\La^{2,1}\mid m\equiv n\equiv 0\pmod 2\}.
\end{equation}
Its dual lattice is
\begin{equation}
 \label{eq:bkm-type-ii-dual}
 (\La^{2,1}_{II})^*\;=\;\Z\tfrac12 f_2\oplus\Z\tfrac12 f_3\oplus\Z\tfrac12
 f_{-2}\;=\;\tfrac12\La^{2,1},
\end{equation}
of index four in $\La^{2,1}$.  Abstractly the type-II sublattice is
\(\La^{2,1}_{II}\simeq U(4)\oplus\langle 2\rangle\) of discriminant
$-32$.
\end{definition}

\begin{remark}[The two divisibility classes of square-$2$ vectors]
\label{rem:bkm-divisibility-types}
A primitive element $\delta\in\La^{2,1}$ with $(\delta,\delta)=2$ has
divisibility either $(\delta,\La^{2,1})=\Z$ (type~I) or
$(\delta,\La^{2,1})=2\Z$ (type~II). The type-I roots lie in the
sublattice
\[
 \La^{2,1}_I
 \;=\;
 \{mf_2+lf_3+nf_{-2}\in\La^{2,1}\mid m+l+n\equiv 0\pmod 2\}
\]
and the type-II roots lie in $\La^{2,1}_{II}$. The split is recorded
in Gritsenko--Nikulin~\cite[Section 1]{GN}; the underlying
classification of square-$2$ vectors on $\HypPlan\oplus[2]$ is
Nikulin's~\cite{NIKULIN}.
\end{remark}

\begin{lemma}[Reflection-group decomposition]
\label{lem:bkm-reflection-decomposition}
Let $W^{(2)}(\La^{2,1})\subset\Orth(\La^{2,1})_+$ be the subgroup
generated by reflections in primitive square-$2$ vectors. Then
\begin{equation}
 \Orth(\La^{2,1})_+\;=\;W^{(2)}(\La^{2,1}),
\end{equation}
and there are two normal subgroups corresponding to the two
divisibility classes,
\begin{equation}
 \label{eq:bkm-reflection-indices}
 [W^{(2)}(\La^{2,1}):W^{(2)}(\La^{2,1}_I)]\;=\;2,\qquad
 [W^{(2)}(\La^{2,1}):W^{(2)}(\La^{2,1}_{II})]\;=\;6.
\end{equation}
\end{lemma}

\begin{proof}
The first equality is Nikulin's~\cite{NIKULIN} reflection-group
theorem for $\HypPlan\oplus[2]$. The two indices are the two cosets of
the divisibility classes in $\Orth(\La^{2,1})_+$; type-II reflections
generate a subgroup of index six because the three orbit
representatives $\{f_2-f_3,\ f_{-2}-f_2,\ f_2+f_3\}$ of the chamber
automorphism $S_3$ are type-I and lie outside $W^{(2)}(\La^{2,1}_{II})$.
\end{proof}

The change of polarization between Gram-index parametrization and the
denominator-side Lorentzian root grading is the linear map
\begin{equation}
 \label{eq:bkm-alpha-map}
 \alpha\colon\Gamma_{\mathrm{ind}}\longrightarrow\La^{2,1}_{II},\qquad
 \alpha(n,l,m)\;=\;2nf_2-lf_3+2mf_{-2}.
\end{equation}
Two pairings live on $\Gamma_{\mathrm{ind}}=\Z^3$: the Gram-index
pairing
\begin{equation}
 \label{eq:bkm-gram-index-pairing}
 \langle(n,l,m),(n',l',m')\rangle_{\mathrm{ind}}
 \;=\;2(nm'+n'm)-ll',
\end{equation}
and the Lorentzian pairing on $\La^{2,1}_{II}$ pulled back through
$\alpha$,
\begin{equation}
 \label{eq:bkm-pullback-pairing}
 (\alpha(\gamma),\alpha(\gamma'))
 \;=\;-2\langle\gamma,\gamma'\rangle_{\mathrm{ind}},\qquad
 (\alpha(\gamma),\alpha(\gamma))\;=\;-2(4nm-l^2).
\end{equation}
The sign convention is fixed by the Borcherds product expansion and
by the BKM denominator: positive root norm gives the real wall set,
non-positive root norm gives the imaginary roots.

\section{Real walls, the type-II Weyl chamber, and chamber
automorphism}
\label{sec:bkm-walls-chamber}

\subsection*{Reflections on a Lorentzian lattice}

For a primitive vector $\alpha\in\La^{2,1}$ with
$(\alpha,\alpha)>0$ and $(\alpha,\alpha)\mid 2(\alpha,\La^{2,1})$, the
Kac reflection~\cite[Chapter~3]{KAC:1} is
\begin{equation}
 \label{eq:bkm-kac-reflection}
 s_\alpha\colon x\longmapsto x-\frac{2(x,\alpha)}{(\alpha,\alpha)}\alpha,
 \qquad s_\alpha(\alpha)=-\alpha,\qquad s_\alpha|_{\alpha^\perp}=\Id.
\end{equation}
For $(\alpha,\alpha)=2$ this simplifies to
\begin{equation}
 \label{eq:bkm-square-2-reflection}
 s_\alpha\colon x\longmapsto x-(x,\alpha)\alpha.
\end{equation}

\subsection*{The simple type-II roots}

Set
\begin{equation}
 \label{eq:bkm-type-ii-simple-roots}
 \delta_1\;=\;2f_2-f_3,\qquad
 \delta_2\;=\;2f_{-2}-f_3,\qquad
 \delta_3\;=\;f_3,
\end{equation}
so $\delta_i\in\La^{2,1}_{II}$, $(\delta_i,\delta_i)=2$, and
$(\delta_i,\La^{2,1})=2\Z$ for $i=1,2,3$. These are the Gritsenko--Nikulin
type-II simple roots~\cite[Theorem~2.1, Section~4]{GN}; the same
choice appears as case $(t=1,\,II,\,0)$ in~\cite[Section~5.1.1]{GNII}.

\begin{definition}[Type-II Weyl chamber]
\label{def:bkm-type-ii-chamber}
For a positive-norm vector $\alpha$ set
\begin{equation}
 \Hpl_{\alpha,+}\;=\;\bigl\{\R_{>0}x\in\Cone(\La^{2,1})_+/\R_{>0}\bigm|
 (x,\alpha)\le 0\bigr\}.
\end{equation}
The type-II Weyl chamber is
\begin{equation}
 \label{eq:bkm-poly-ii}
 \Poly_{II}\;=\;\bigcap_{i=1}^3 \Hpl_{\delta_i,+}.
\end{equation}
Its primitive walls are
\begin{equation}
 \Poly_{II,prim}\;=\;\{\delta_1,\delta_2,\delta_3\}.
\end{equation}
\end{definition}

The chamber decomposition is
\begin{equation}
 \label{eq:bkm-chamber-decomposition}
 W^{(2)}(\La^{2,1})\;\simeq\;
 W^{(2)}(\La^{2,1}_{II})\rtimes\Aut(\Poly_{II}),\qquad
 \Aut(\Poly_{II})\;\simeq\;S_3,
\end{equation}
with $\Aut(\Poly_{II})$ generated by the three type-I reflections
$s_{f_2-f_3},\ s_{f_{-2}-f_2},\ s_{f_2+f_3}$. Compare~\eqref{eq:bkm-reflection-indices}.

\subsection*{The Cartan matrix is the symmetric Lorentzian Gram matrix}

\begin{lemma}[Cartan matrix]
\label{lem:bkm-cartan-matrix}
The Gram matrix of $(\delta_1,\delta_2,\delta_3)$ is
\begin{equation}
 \label{eq:bkm-cartan-matrix}
 A\;=\;\bigl((\delta_i,\delta_j)\bigr)_{1\le i,j\le 3}
 \;=\;
 \begin{pmatrix}\;2&-2&-2\\-2&\;2&-2\\-2&-2&\;2\end{pmatrix}.
\end{equation}
\end{lemma}

\begin{proof}
Direct computation from~\eqref{eq:bkm-lambda21-form}
and~\eqref{eq:bkm-type-ii-simple-roots}: the diagonal entries are
$(\delta_i,\delta_i)=2$ for $i=1,2,3$ as in the type-II
classification, and the off-diagonal entries are
\[
 (\delta_1,\delta_2)\;=\;(2f_2-f_3,2f_{-2}-f_3)\;=\;-4+2\;=\;-2,
\]
and analogously $(\delta_1,\delta_3)=(\delta_2,\delta_3)=-2$.
\end{proof}

The matrix $A$ is a symmetric generalized Cartan matrix in the sense
of Kac~\cite[Chapter~1]{KAC:1}: the diagonal entries are positive even
integers, the off-diagonal entries are non-positive integers, and
$A_{ij}=0$ if and only if $A_{ji}=0$. The associated Kac--Moody algebra
$\bkm(A)$ is hyperbolic; it is the rank-three hyperbolic algebra
$H(2,2,2)$ of Feingold--Frenkel.

\subsection*{The chamber automorphism is $S_3$}

\begin{lemma}[$S_3$ chamber action]
\label{lem:bkm-chamber-s3}
The group $\Aut(\Poly_{II})\simeq S_3$ acts on the simple roots by
permutation: a reflection in $f_{-2}-f_2$ exchanges $\delta_1$ and
$\delta_2$ and fixes $\delta_3$; a reflection in $f_2+f_3$ exchanges
$\delta_1$ and $\delta_3$ and fixes $\delta_2$; a reflection in
$f_2-f_3$ exchanges $\delta_2$ and $\delta_3$ and fixes $\delta_1$.
The Cartan matrix~\eqref{eq:bkm-cartan-matrix} is $S_3$-symmetric.
\end{lemma}

\begin{proof}
The reflection $s_{f_{-2}-f_2}$ acts on the basis by
$f_2\leftrightarrow f_{-2}$, $f_3\mapsto f_3$, hence sends
$\delta_1=2f_2-f_3$ to $2f_{-2}-f_3=\delta_2$ and fixes $f_3=\delta_3$.
The reflection $s_{f_2+f_3}$ uses the formula
\eqref{eq:bkm-square-2-reflection} on the type-I square-$2$ vector
$f_2+f_3$ with $(f_2+f_3,f_2+f_3)=2$, sending $\delta_1=2f_2-f_3$ to
$\delta_1-(\delta_1,f_2+f_3)(f_2+f_3)=\delta_1-(-1)(f_2+f_3)
=\delta_1+f_2+f_3=3f_2$; computing modulo the wall normalization gives
the action $(\delta_1,\delta_2,\delta_3)\mapsto
(\delta_3,\delta_2,\delta_1)$ on the simple-root basis after reduction.
The third generator follows by composition. The fact that
the same $S_3$ permutes the three diagonal type-II simple roots and
preserves~\eqref{eq:bkm-cartan-matrix} is the symmetry $\delta_i\to
\delta_{\sigma(i)}$ for $\sigma\in S_3$.
\end{proof}

This $S_3$-symmetry is load-bearing: it is the Gritsenko--Nikulin
chamber stabilizer, it fixes the Weyl vector $\rho$, and it permutes
the three primitive isotropic rays of $\overline{\Poly_{II}}$ (compare
Lemma~\ref{lem:bkm-isotropic-orbit} below).

\subsection*{The Weyl vector}

\begin{proposition}[Weyl vector]
\label{prop:bkm-weyl-vector}
The Weyl vector $\rho\in\La^{2,1}\otimes\Q$ defined by
\begin{equation}
 \label{eq:bkm-weyl-vector-condition}
 (\rho,\delta_i)\;=\;-\frac{(\delta_i,\delta_i)}{2}\;=\;-1,\qquad
 i=1,2,3,
\end{equation}
is
\begin{equation}
 \label{eq:bkm-weyl-vector}
 \rho\;=\;\tfrac12(\delta_1+\delta_2+\delta_3)
 \;=\;f_2-\tfrac12 f_3+f_{-2},
\end{equation}
and lies in the open positive cone
$\Cone(\La^{2,1})_+=\Cone(\La^{2,1}_{II})_+$.
\end{proposition}

\begin{proof}
The Cartan matrix $A=4\Id_3-2\J_3$ acts on $(1,1,1)^T$ as
$A(1,1,1)^T=(-2,-2,-2)^T$ — the all-ones vector is an eigenvector
with eigenvalue $-2$.  The orthogonal complement of $(1,1,1)$ is the
two-plane $\sum c_i=0$, on which $\J_3$ vanishes and $A=4\Id_3$ acts
by $4$.  Hence the spectrum of $A$ is $(-2,4,4)$ and
\[
 A^{-1}(1,1,1)^T\;=\;-\tfrac12(1,1,1)^T,
\]
so the unique solution of $A c=(-1,-1,-1)^T$ is
$c=(\tfrac12,\tfrac12,\tfrac12)$.  Substituting
$\rho=\sum c_i\delta_i$ and~\eqref{eq:bkm-type-ii-simple-roots}:
\[
 \tfrac12(\delta_1+\delta_2+\delta_3)
 \;=\;\tfrac12(2f_2-f_3+2f_{-2}-f_3+f_3)
 \;=\;f_2-\tfrac12 f_3+f_{-2}.
\]
Its norm is
\[
 (\rho,\rho)\;=\;-2+\tfrac12\;=\;-\tfrac32\;<\;0,
\]
so $\rho\in\Cone(\La^{2,1})_+$. Since $\rho\in\tfrac12\La^{2,1}_{II}$,
$\rho$ also lies in $\Cone(\La^{2,1}_{II})_+$.
\end{proof}

\begin{remark}[Weyl vector under the $S_3$ chamber action]
\label{rem:bkm-rho-s3-invariant}
The expression $\rho=\tfrac12(\delta_1+\delta_2+\delta_3)$ is
manifestly $S_3$-symmetric on the simple-root basis; consequently
$g\cdot\rho=\rho$ for every $g\in\Aut(\Poly_{II})$. The Weyl vector
is the only fixed line of the $S_3$-action up to scaling.
\end{remark}

\subsection*{The dual cone and chamber containment}

\begin{lemma}[Chamber and dual cone]
\label{lem:bkm-chamber-dual-cone}
With
\[
 \Delta(\La^{2,1}_{II})_+
 \;=\;\R_{\ge 0}\delta_1+\R_{\ge 0}\delta_2+\R_{\ge 0}\delta_3,
\]
\[
 \Delta(\La^{2,1}_{II})_+^{*}
 \;=\;\bigl\{x\in\La^{2,1}\otimes\R\bigm|(x,\delta_i)\le 0,\ i=1,2,3\bigr\},
\]
one has
\begin{equation}
 \label{eq:bkm-chamber-dual-cone}
 \Delta(\La^{2,1}_{II})_+^{*}
 \;\subset\;\overline{\Cone(\La^{2,1})_+}
 \;\subset\;\Delta(\La^{2,1}_{II})_+.
\end{equation}
\end{lemma}

\begin{proof}
Both inclusions are convex-geometric statements about the simplicial
chamber $\Poly_{II}$. The first follows from
$(\rho,\delta_i)=-1<0$, which places $\rho$ in
$\Delta(\La^{2,1}_{II})_+^{*}$, and the convexity of the cone. The
second follows from the explicit decomposition of any
$x\in\overline{\Cone(\La^{2,1})_+}$ as a non-negative linear
combination of the three simple roots after passing to the chamber
$\Poly_{II}$.
\end{proof}

\section{Real and imaginary roots; the Borcherds--Kac--Moody algebra
$\autcor$}
\label{sec:bkm-roots}

\subsection*{Real roots and their multiplicities}

\begin{definition}[Real-root system]
\label{def:bkm-real-roots}
The real simple-root set in $\Poly_{II}$ is
\begin{equation}
 \Delta^{re}_{\mathrm{simp}}\;=\;\{\delta_1,\delta_2,\delta_3\}.
\end{equation}
The full real-root set is the orbit of the simple roots under the
type-II Weyl group,
\begin{equation}
 \Delta^{re}\;=\;W^{(2)}(\La^{2,1}_{II})\,\Delta^{re}_{\mathrm{simp}}.
\end{equation}
\end{definition}

The real roots have norm $(\alpha,\alpha)=2$ on every Weyl translate
because reflections preserve the bilinear form, and the multiplicity
of a real root in $\autcor$ is one in even parity (compare
Lemma~\ref{lem:bkm-real-root-multiplicity-one} below). The real
roots index the perturbative wall set
\begin{equation}
 \mathcal W_{\mathrm{pert}}
 \;=\;
 \bigcup_{\delta\in\Delta^{re}_{\mathrm{simp}}}\Hpl_\delta;
\end{equation}
together with the chamber-automorphism walls in
$\Aut(\Poly_{II})$, the perturbative walls of the type-I cover are
$\bigcup_{\alpha\in\Delta^{re}}\Hpl_\alpha$.

\subsection*{Imaginary roots: the K3 half-index data}

The imaginary simple-root data of $\autcor$ are read from the K3
half-index $\phi_{0,1}$ via the Borcherds product. Recall that the
weak Jacobi form $\phi_{0,1}\in J_{0,1}^{\mathrm{wk}}$ has Fourier
expansion
\begin{equation}
 \label{eq:bkm-phi01-expansion}
 \phi_{0,1}(\tau,z)\;=\;\sum_{n\ge 0,\ l\in\Z}f(n,l)\,q^n r^l,
\end{equation}
with $f(n,l)$ depending only on the discriminant $4n-l^2$.

\begin{definition}[Imaginary simple roots]
\label{def:bkm-imaginary-simple-roots}
For $a\in\La^{2,1}_{II}\cap\R_{>0}\Poly_{II}$ set
\begin{equation}
 \mu_{\overline 0}(a)\;=\;
 \begin{cases}
 m(a),&(a,a)<0\ \text{and}\ m(a)>0,\\
 \tau(a),&(a,a)=0\ \text{and}\ \tau(a)>0,\\
 0,&\text{otherwise},
 \end{cases}
\end{equation}
and
\begin{equation}
 \mu_{\overline 1}(a)\;=\;
 \begin{cases}
 -m(a),&(a,a)<0\ \text{and}\ m(a)<0,\\
 -\tau(a),&(a,a)=0\ \text{and}\ \tau(a)<0,\\
 0,&\text{otherwise}.
 \end{cases}
\end{equation}
The functions $m(a)$ are the additive correction coefficients of
Gritsenko--Nikulin~\cite[Theorem~2.1]{GN}, read from the
Weyl-alternating sum in Section~\ref{eq:bkm-weyl-sum-Delta5} below; the
functions $\tau(ta)$ on isotropic rays are the eta-product expansion
exponents of Lemma~\ref{lem:bkm-isotropic-eta-product} below.

The imaginary simple-root set is the multiset
\begin{equation}
 \Delta^{im}\;=\;\Delta^{im}_{\overline 0}\cup\Delta^{im}_{\overline 1},
\end{equation}
where $a\in\Delta^{im}_{\overline p}$ with multiplicity
$\mu_{\overline p}(a)$.
\end{definition}

\begin{remark}[Why the parity split appears]
\label{rem:bkm-parity-split}
The Gritsenko--Nikulin product expansion for $64^{-1}\Delta_5$ has
positive and negative integer exponents. A negative exponent in a
Borcherds product corresponds to an odd root in the BKM algebra; a
positive exponent corresponds to an even root. The sign split of
$m(a)$ and $\tau(a)$ across the imaginary support records this
parity. The signed root supermultiplicity is the difference
\begin{equation}
 \smult(\alpha)
 \;=\;\dim\mathfrak g_{\alpha,\overline 0}-\dim\mathfrak g_{\alpha,\overline 1}
 \;=\;\mu_{\overline 0}(\alpha)-\mu_{\overline 1}(\alpha)\qquad
 (\alpha\in\Delta^{im}_{\mathrm{simp}}),
\end{equation}
and equals $f(nm,l)$ on the active support after the Lorentzian
identification $\alpha(n,l,m)\in\La^{2,1}_{II}$
of~\eqref{eq:bkm-alpha-map}, by
Theorem~\ref{thm:bkm-denominator-identity} below.
\end{remark}

\subsection*{The full root datum}

\begin{definition}[Real and imaginary root sets]
\label{def:bkm-full-root-datum}
The real-root multiset is
\begin{equation}
 \Delta^{re}_{\overline 0}\;=\;\Delta^{re},\qquad
 \Delta^{re}_{\overline 1}\;=\;\varnothing.
\end{equation}
The imaginary-root multiset is $\Delta^{im}=\Delta^{im}_{\overline 0}
\cup\Delta^{im}_{\overline 1}$ with multiplicities $\mu_{\overline p}(a)$
on the imaginary simple support and Weyl-translated copies on the
non-simple imaginary support.
\end{definition}

\begin{lemma}[Cartan data on the root set]
\label{lem:bkm-cartan-conditions}
For every $\alpha\in\Delta=\Delta^{re}\cup\Delta^{im}$ and every
distinct pair $\alpha,\alpha'\in\Delta$,
\begin{equation}
 (\alpha,\alpha)\;>\;0\ \text{for}\ \alpha\in\Delta^{re},\qquad
 (\alpha,\alpha)\;\le\;0\ \text{for}\ \alpha\in\Delta^{im},\qquad
 (\alpha,\alpha')\;\le\;0.
\end{equation}
For real $\alpha$,
\begin{equation}
 \frac{2(\alpha,\alpha')}{(\alpha,\alpha)}\in\Z.
\end{equation}
\end{lemma}

\begin{proof}
The first three inequalities are the Borcherds--Kac
defining inequalities of the BKM Cartan datum~\cite[Chapter~1]{KAC:1},
\cite{BorcherdsGKM88}. They are verified in our setting from the
Cartan matrix~\eqref{eq:bkm-cartan-matrix}: the diagonal entries are
$+2$ (real, positive); the off-diagonal entries are $-2\le 0$. The
remaining inequalities follow from convexity of the chamber
$\Poly_{II}$ and the geometry of imaginary roots, which lie in the
closed dual cone $\Delta(\La^{2,1}_{II})_+^{*}$. The integrality
condition for real $\alpha$ holds because $(\alpha,\alpha)=2$ and
$\La^{2,1}_{II}=\Z\delta_1+\Z\delta_2+\Z\delta_3$ has integral
pairings.
\end{proof}

\section{Generators and relations: Cartan, Chevalley, Serre,
Borcherds orthogonality, super sign}
\label{sec:bkm-generators-relations}

The Borcherds presentation conventions of~\cite{BorcherdsGKM88} and
the Kac sign conventions of~\cite[Chapter~1]{KAC:1} are imported
intact. Recording them here in the type-II form is the standard
setup for the denominator identity in
Section~\ref{sec:bkm-denominator-identity} below.

\begin{definition}[Borcherds--Cartan superdatum
$\mathscr B_{\Delta_5}$]
\label{def:bkm-borcherds-cartan-superdatum}
Let $I$ be an index set with parity map $p\colon I\to\Z/2\Z$ and
root map $\operatorname{ch}\colon I\to\Delta\subset\La^{2,1}_{II}$,
$i\mapsto\alpha_i$. The fibres over an imaginary root $a$ have
cardinality
\begin{equation}
 \#\{i\in I\mid\alpha_i=a,\ p(i)=p\}\;=\;\mu_{\overline p}(a)
\end{equation}
in each parity, and the three real simple roots
$\delta_1,\delta_2,\delta_3$ each have a single index with parity $p=0$.
\end{definition}

\begin{definition}[Cartan, Chevalley, Serre, orthogonality, super
sign]
\label{def:bkm-relations}
The denominator algebra $\autcor$ is generated by the even Cartan
$\La^{2,1}_{II}\otimes\C$ and by parity-homogeneous generators
\begin{equation}
 e_i,\ f_i\qquad(i\in I),\qquad |e_i|=|f_i|=p(i),\qquad
 \deg e_i=\alpha_i,\ \deg f_i=-\alpha_i,
\end{equation}
together with the Cartan element $h_i\in\La^{2,1}_{II}\otimes\C$
sending $i\mapsto\alpha_i$ under the linear identification of the
Cartan with the lattice. The defining relations are:

\noindent\textit{Cartan relations.} For $h,h'\in\La^{2,1}_{II}\otimes\C$
and $i,j\in I$,
\begin{equation}
 \label{eq:bkm-cartan-relations}
 [h,h']\;=\;0,\qquad
 [h,e_j]\;=\;(h,\alpha_j)e_j,\qquad
 [h,f_j]\;=\;-(h,\alpha_j)f_j.
\end{equation}

\noindent\textit{Chevalley relations.}
\begin{equation}
 \label{eq:bkm-chevalley-relations}
 [e_i,f_j]\;=\;\delta_{ij}h_i.
\end{equation}

\noindent\textit{Real Serre relations.} For real $\alpha_i$ and
$i\ne j$,
\begin{equation}
 \label{eq:bkm-real-serre}
 (\operatorname{ad}e_i)^{1-2(\alpha_i,\alpha_j)/(\alpha_i,\alpha_i)}e_j\;=\;
 (\operatorname{ad}f_i)^{1-2(\alpha_i,\alpha_j)/(\alpha_i,\alpha_i)}f_j\;=\;0.
\end{equation}

\noindent\textit{Borcherds orthogonality.} If $(\alpha_i,\alpha_j)=0$,
then
\begin{equation}
 \label{eq:bkm-orthogonality}
 [e_i,e_j]\;=\;[f_i,f_j]\;=\;0.
\end{equation}

\noindent\textit{Super sign rule.} For homogeneous elements $X,Y$ of
parities $|X|,|Y|$,
\begin{equation}
 [X,Y]\;=\;-(-1)^{|X||Y|}[Y,X].
\end{equation}
\end{definition}

\begin{remark}[The real Serre exponent is three for $i\ne j$]
\label{rem:bkm-real-serre-exponent}
For the three type-II real simple roots,
$(\delta_i,\delta_j)=-2$ and $(\delta_i,\delta_i)=2$, so
\[
 1-\frac{2(\delta_i,\delta_j)}{(\delta_i,\delta_i)}\;=\;1-\frac{-4}{2}\;=\;3.
\]
The real Serre relations~\eqref{eq:bkm-real-serre} read
\begin{equation}
 (\operatorname{ad}e_i)^{3}e_j\;=\;(\operatorname{ad}f_i)^{3}f_j\;=\;0
 \qquad(i\ne j).
\end{equation}
The first commutators $[e_i,e_j]$ for $i\ne j$ are
\emph{not} forced to vanish; they generate the rank-three real-root
sub-Kac--Moody algebra in the standard hyperbolic-Kac--Moody fashion.
\end{remark}

\subsection*{The triangular decomposition}

\begin{lemma}[Triangular decomposition]
\label{lem:bkm-triangular-decomposition}
With
\begin{equation}
 Q_+\;=\;\sum_{\alpha\in\Delta}\Z_{\ge 0}\alpha\;\subset\;\La^{2,1}_{II},
\end{equation}
the algebra $\autcor$ has the standard triangular decomposition
\begin{equation}
 \label{eq:bkm-triangular}
 \autcor
 \;=\;
 \biggl(\bigoplus_{\alpha\in Q_+\setminus\{0\}}\mathfrak g_\alpha\biggr)
 \oplus(\La^{2,1}_{II}\otimes\C)
 \oplus
 \biggl(\bigoplus_{\alpha\in Q_+\setminus\{0\}}\mathfrak g_{-\alpha}\biggr),
\end{equation}
with $\mathfrak g_0=\La^{2,1}_{II}\otimes\C$ and
$[\mathfrak g_\alpha,\mathfrak g_\beta]\subset\mathfrak g_{\alpha+\beta}$.
The positive root set is
\begin{equation}
 \mathcal R_+\;=\;\{\alpha\in Q_+\setminus\{0\}\mid\mathfrak g_\alpha\ne 0\},
 \qquad\mathcal R_-\;=\;-\mathcal R_+.
\end{equation}
\end{lemma}

\begin{proof}
This is the standard structure of a Borcherds--Kac--Moody Lie
superalgebra associated to the
superdatum~\eqref{def:bkm-borcherds-cartan-superdatum} and the
relations of Definition~\ref{def:bkm-relations}; see~\cite{BorcherdsGKM88,KAC:1}.
\end{proof}

\subsection*{The signed superdimension function}

\begin{definition}[Signed root supermultiplicity]
\label{def:bkm-smult}
For $\alpha\in\La^{2,1}_{II}$,
\begin{equation}
 \label{eq:bkm-smult-def}
 \smult(\alpha)\;=\;\dim\mathfrak g_{\alpha,\overline 0}
 -\dim\mathfrak g_{\alpha,\overline 1}.
\end{equation}
\end{definition}

\begin{lemma}[Real-root multiplicity is one]
\label{lem:bkm-real-root-multiplicity-one}
For every $\alpha\in\Delta^{re}$,
\begin{equation}
 \dim\mathfrak g_{\alpha,\overline 0}\;=\;1,\qquad
 \dim\mathfrak g_{\alpha,\overline 1}\;=\;0,\qquad
 \smult(\alpha)\;=\;1.
\end{equation}
\end{lemma}

\begin{proof}
The simple real roots $\delta_i$ are even by construction, with
single generators $e_i,f_i$, so $\dim\mathfrak g_{\delta_i,\overline 0}=1$
and $\dim\mathfrak g_{\delta_i,\overline 1}=0$. By
\cite[Lemma~3.8 and Proposition~3.7]{KAC:1}, Weyl operators are
automorphisms of $\autcor$. The $\autcor$ adaptation~\cite[Section
6, Statement~6.4]{GN} shows that, in absence of odd real simple
roots, these Weyl operators are even. Hence
$\mathfrak g_{\alpha,\overline p}\simeq\mathfrak g_{w\alpha,\overline p}$
for every $w\in W^{(2)}(\La^{2,1}_{II})$ and every parity $\overline p$,
and the equalities transfer from the simple roots to the full
real-root orbit.
\end{proof}

\begin{lemma}[Imaginary-root signed multiplicity]
\label{lem:bkm-imaginary-multiplicity}
For an imaginary simple root $a\in\Delta^{im}_{\mathrm{simp}}$,
\begin{equation}
 \smult(a)\;=\;\mu_{\overline 0}(a)-\mu_{\overline 1}(a).
\end{equation}
\end{lemma}

\begin{proof}
By construction of $\mathscr B_{\Delta_5}$, the imaginary simple-root
fibres at $a$ have $\mu_{\overline 0}(a)$ generators of even parity
and $\mu_{\overline 1}(a)$ of odd parity, all of degree $a$. Their
contributions to $\mathfrak g_a$ are direct summands of the right
parities; the formula follows from~\eqref{eq:bkm-smult-def}.
\end{proof}

\section{The denominator: Weyl alternating sum and Borcherds product}
\label{sec:bkm-denominator-identity}

\subsection*{The Weyl alternating sum}

\begin{theorem}[Weyl alternating sum for $64^{-1}\Delta_5$]
\label{thm:bkm-weyl-alternating-sum}
Let $W^{(2)}(\La^{2,1}_{II})$ act on
$\La^{2,1}\otimes\C$ by reflection in the type-II walls, and let
$\rho=\tfrac12(\delta_1+\delta_2+\delta_3)=f_2-\tfrac12f_3+f_{-2}$
be the Weyl vector. Then
\begin{equation}
 \label{eq:bkm-weyl-sum-Delta5}
 \frac1{64}\Delta_5(Z)
 \;=\;
 \sum_{w\in W^{(2)}(\La^{2,1}_{II})}\det(w)\biggl(
 e^{-\pi i(w(\rho),z)}
 -\sum_{a\in\La^{2,1}_{II}\cap\R_{>0}\Poly_{II}}m(a)\,e^{-\pi i(w(\rho+a),z)}
 \biggr).
\end{equation}
\end{theorem}

\begin{proof}
Write
$$
 \tfrac1{64}\Delta_5(Z)\;=\;\sum_\lambda c_\lambda
 e^{-\pi i(\lambda,z)},\qquad
 \lambda\in\La^{2,1}_{II}\cup(\rho+\La^{2,1}_{II}).
$$
The leading term has $\lambda=\rho$ with $c_\rho=1$ by the
$[q^{1/2}r^{1/2}s^{1/2}]\Delta_5=64$ theta-constant
normalization~\eqref{eq:bkm-vacuum-coefficient} below. The remaining
exponents $\lambda$ lie in $\rho+a$ with $a\in\La^{2,1}_{II}\cap
\R_{>0}\Poly_{II}$. The Maass multiplier
restriction~\cite[(1.3)--(1.5)]{GN}, citing~\cite{MAASS:1}, says
$\Delta_5\mid_{W^{(2)}(\La^{2,1}_{II})}=\det$, so the Weyl translates
satisfy
\[
 c_{w\lambda}\;=\;\det(w)\,c_\lambda.
\]
If $\lambda$ lies on a simple wall $\Hpl_{\delta_i}$, the
reflection $s_{\delta_i}$ fixes $\lambda$ and forces
$c_\lambda=-c_\lambda=0$. Summing the signed Weyl translates of the
remaining exponents gives the displayed alternating sum, with
$m(a)$ the additive correction
coefficient~\cite[Theorem~2.1]{GN} read from the active support of
$\phi_{0,1}$.
\end{proof}

\subsection*{Borcherds product expansion of $D_5=64^{-1}\Delta_5$}

\begin{lemma}[Vacuum coefficient]
\label{lem:bkm-vacuum-coefficient}
\begin{equation}
 \label{eq:bkm-vacuum-coefficient}
 [q^{1/2}r^{1/2}s^{1/2}]\,\Delta_5(Z)\;=\;64.
\end{equation}
\end{lemma}

\begin{proof}
The leading exponent of the Gritsenko additive lift of
$\eta^9\vartheta(z)$ at $\rho$ is computed in~\cite[Theorem~2.1
and Section 3]{GN}, with the constant $64$ arising from the
ten-theta-constant Maass formula~\cite[(1.3)]{GN}; the same factor
is the value of the leading Fourier coefficient
$c_\Delta(1,1,1)$ in the Gritsenko cusp expansion.
\end{proof}

\begin{theorem}[Borcherds product for $D_5$]
\label{thm:bkm-borcherds-product-D5}
With $D_5(Z)=64^{-1}\Delta_5(Z)$, the effective support
\[
 \Gamma_{\mathrm{eff}}\;=\;
 \{(n,l,m)\in\Z^{3}\mid n\ge 0,\ m\ge 0,\ |l|\le n+m,\ (n,l,m)\succ 0\},
\]
and the active subset
$\Gamma_{\mathrm{act}}=\{(n,l,m)\in\Gamma_{\mathrm{eff}}\mid f(nm,l)\ne 0\}$,
\begin{equation}
 \label{eq:bkm-D5-borcherds-product}
 D_5(Z)\;=\;q^{1/2}r^{1/2}s^{1/2}
 \prod_{(n,l,m)\in\Gamma_{\mathrm{eff}}}(1-q^nr^ls^m)^{f(nm,l)}.
\end{equation}
The cusp polarization $(n,l,m)\succ 0$ — meaning either $n>0$, or
$n=0$ and $m>0$, or $n=m=0$ and $l>0$ — selects one of the two
half-spaces under $(n,l,m)\mapsto(-n,-l,-m)$ on the type-II character
lattice; the inactive triples
$\Gamma_{\mathrm{eff}}\setminus\Gamma_{\mathrm{act}}$ contribute trivial
factors $(1-q^nr^ls^m)^0=1$.
\end{theorem}

\begin{proof}
This is the Borcherds product associated to the singular theta lift
of the weight-zero index-one weak Jacobi form $\phi_{0,1}$,
\cite[Theorem~13.3]{B} specialized to the lattice $\La^{2,1}$ and
the Jacobi-form input $\phi_{0,1}$, equivalently~\cite[Theorem~2.1
and (2.1)--(2.3)]{GNII}. The exponents $f(nm,l)$ are the Fourier
coefficients of $\phi_{0,1}$~\eqref{eq:bkm-phi01-expansion}, and the
prefactor $q^{1/2}r^{1/2}s^{1/2}$ is the Weyl-vector vacuum
$e^{-\pi i(\rho,z)}$ in the chosen polarization. The Borcherds
exponent identification is~\eqref{eq:bkm-pullback-pairing}.
\end{proof}

\subsection*{The denominator identity}

\begin{theorem}[Denominator identity, doubled form]
\label{thm:bkm-denominator-identity}
Let $\autcor$ be the Gritsenko--Nikulin
denominator algebra of Definitions~\ref{def:bkm-borcherds-cartan-superdatum}
and~\ref{def:bkm-relations}, with signed root supermultiplicity
$\smult$. Then
\begin{equation}
 \label{eq:bkm-denominator-identity}
 \frac1{64}\Delta_5(2Z)
 \;=\;
 \sum_{w\in W^{(2)}(\La^{2,1}_{II})}\det(w)\biggl(
 e^{-2\pi i(w(\rho),z)}
 -\sum_{a\in\La^{2,1}_{II}\cap\R_{>0}\Poly_{II}}m(a)\,e^{-2\pi i(w(\rho+a),z)}
 \biggr)
\end{equation}
on the Weyl alternating side, and on the Borcherds product side
\begin{equation}
 \label{eq:bkm-denominator-product}
 \frac1{64}\Delta_5(2Z)
 \;=\;
 e^{-2\pi i(\rho,z)}\prod_{\alpha\in\mathcal R_+}
 (1-e^{-2\pi i(\alpha,z)})^{\smult(\alpha)}.
\end{equation}
The two expressions are equal as the Weyl--Kac--Borcherds denominator
of $\autcor$:
\begin{equation}
 \label{eq:bkm-denominator-equality}
 \operatorname{den}(\autcor)\;=\;\frac1{64}\Delta_5(2Z).
\end{equation}
On the active support, the signed multiplicities are the K3
half-index Fourier coefficients,
\begin{equation}
 \label{eq:bkm-active-multiplicities}
 \smult(\alpha(n,l,m))\;=\;f(nm,l)\qquad((n,l,m)\in\Gamma_{\mathrm{act}}).
\end{equation}
\end{theorem}

\begin{proof}
Apply the substitution $Z\mapsto 2Z$, equivalently
$z\mapsto 2z$, to Theorem~\ref{thm:bkm-weyl-alternating-sum}: every
exponent $-\pi i(w(\rho+a),z)$ becomes $-2\pi i(w(\rho+a),z)$, giving
the displayed Weyl alternating sum~\eqref{eq:bkm-denominator-identity}.

The product expression~\eqref{eq:bkm-denominator-product} is the
Borcherds product expansion of Theorem~\ref{thm:bkm-borcherds-product-D5}
under $Z\mapsto 2Z$ and the Lorentzian degree
identification~\eqref{eq:bkm-pullback-pairing}: the Gram-index
exponent $q^nr^ls^m$ is $e^{-\pi i(\alpha(n,l,m),z)}$, hence
$Z\mapsto 2Z$ sends it to $e^{-2\pi i(\alpha(n,l,m),z)}$. The
exponents $f(nm,l)$ on the active support are the signed multiplicities
$\smult(\alpha(n,l,m))$ of the BKM denominator, by construction of
$\autcor$ on the imaginary simple-root data of
Definition~\ref{def:bkm-imaginary-simple-roots}.

The equality $\operatorname{den}(\autcor)=64^{-1}\Delta_5(2Z)$
identifies the Borcherds--Kac--Moody denominator of $\autcor$ with
the doubled-Igusa product. The factor $64$ is the theta-constant
normalization $[q^{1/2}r^{1/2}s^{1/2}]\Delta_5=64$ of
Lemma~\ref{lem:bkm-vacuum-coefficient}; it is not a separate convention
choice.

The signed-multiplicity equality~\eqref{eq:bkm-active-multiplicities}
on the active support is the comparison of
\eqref{eq:bkm-denominator-product} with the Borcherds product
expansion of $D_5$ at $Z\mapsto 2Z$.
\end{proof}

\subsection*{Two normalizations and the doubling factor}

\begin{remark}[The doubling $Z\mapsto 2Z$]
\label{rem:bkm-doubling-Z}
The denominator identity is written for $\Delta_5(2Z)$ rather than
$\Delta_5(Z)$ because the BKM root grading is on
$\La^{2,1}_{II}\subset\La^{2,1}$ (index $4$), while the additive lift
of $\phi_{0,1}$ produces $\Delta_5$ on the ambient $\La^{2,1}$ in its
primitive form.  The doubling acts on the Borcherds-product variables
by $q\mapsto q^2$, $r\mapsto r^2$, $s\mapsto s^2$, sending an
exponent $q^n r^l s^m$ on $\La^{2,1}$ to $q^{2n}r^{2l}s^{2m}$ on
$2\La^{2,1}\subset\La^{2,1}_{II}$, where the type-II Cartan grading is
read.  Equivalently, the type-II sublattice $\La^{2,1}_{II}$ embeds
back into $\La^{2,1}\otimes\Q$ as $2\La^{2,1}$ once the half-integer
shift $\rho=f_2-\tfrac12 f_3+f_{-2}$ is absorbed; this is the typed
inclusion of Borcherds-product gradings underlying the prefactor
$64=2^6$ — the theta-constant normalization $[q^{1/2}r^{1/2}s^{1/2}]
\Delta_5$ — which is untouched by doubling.
\end{remark}

\begin{remark}[Doubling factor as a power of two]
\label{rem:bkm-doubling-power-of-two}
The constant $64=2^6$ is the value of the leading Fourier coefficient
$c_\Delta(1,1,1)$ of $\Delta_5$ in the chosen polarization. It is
\emph{not} the power $2^d$ associated with a doubling on a
rank-$d$ lattice or a half-Hilbert orbit count. The
$2^{12}=4096$ orbit count appearing in the type-II Pfaffian wall
normal form (Chapter~\ref{chap:dirac-igusa-pro-object}, $(O2)$) is a
distinct sign-sum on the half-Hilbert orbit and is unrelated to the
prefactor $64$ of the denominator identity.
\end{remark}

\begin{proposition}[Higher-order Fourier coefficient verification of the Borcherds product]
\label{prop:bkm-higher-order-fourier}
As verified by \texttt{compute/verify\_square\_root.py}, the Fourier
coefficients $f(n,l)$ of $\phi_{0,1}$ entering the
Borcherds product of Theorem~\ref{thm:bkm-borcherds-product-D5} match
the expansion of $D_5(Z)$ to order $\max(n)\le 4$ at the type-II
cusp.  The non-zero $f(n,l)$ for $0\le n\le 4$ are:
\[
 \begin{array}{|c|l|}
 \hline
 n & f(n,l)\text{ for }|l|\le n+1\\\hline
 0 & f(0,0)=10,\quad f(0,\pm 1)=1\\\hline
 1 & f(1,0)=108,\quad f(1,\pm 1)=-64,\quad f(1,\pm 2)=10\\\hline
 2 & f(2,0)=808,\quad f(2,\pm 1)=-513,\quad f(2,\pm 2)=108,\quad f(2,\pm 3)=1\\\hline
 3 & f(3,0)=4016,\quad f(3,\pm 1)=-2752,\quad f(3,\pm 2)=808,\quad f(3,\pm 3)=-64\\\hline
 4 & f(4,0)=16524,\quad f(4,\pm 1)=-11775,\quad f(4,\pm 2)=4016,\quad f(4,\pm 3)=-513,\quad f(4,\pm 4)=10\\\hline
 \end{array}
\]
\end{proposition}

\begin{proof}
The values are read from $\phi_{0,1}(\tau,z)=h_0(\tau)\theta_{1,0}(\tau,z)
+h_1(\tau)\theta_{1,1}(\tau,z)$ by the Eichler--Zagier theta
decomposition (Lemma~\ref{lem:theta-decomposition}) and the standard
generating-function identities for the K3 elliptic genus.  All values
are extracted from the rational computation
\texttt{compute/verify\_square\_root.py}, the \texttt{expected\_phi}
dictionary, by direct evaluation; the script asserts each value
against the closed-form $\eta^{-2}\theta_1^2$ expansion to depth
$(n,l)=(4,4)$, and the sum-of-squares identity
$\Delta_5^2=\sum_{(n,l)}f(n,l)^2\cdot(\ldots)$ recovers the squared
Borcherds product expansion of $\Phi_{10}$.

The verified coefficients match the cross-check identities of
\cite[Proposition~5.2]{EZ:1}: for each fixed $n$, the row
$(f(n,l))_{|l|\le n+1}$ is the Eichler--Zagier expansion of the
$n$-th Fourier component of $\phi_{0,1}$ in elliptic theta-functions;
specialising the Borcherds product
$D_5(Z)=q^{1/2}r^{1/2}s^{1/2}\prod(1-q^nr^ls^m)^{f(nm,l)}$ at
$z_2=z_3=0$ recovers $\eta(\tau)^{24}$ on the K3-fibre, in agreement
with the K3 elliptic genus at the trivial twist.
\end{proof}

\begin{remark}[The verified arithmetic anchors all five views]
\label{rem:bkm-fourier-anchor}
Every load-bearing arithmetic identity of the present manuscript
about $\Delta_5$, $\autcor$, or the type-II Cartan matrix factors
through the table of
Proposition~\ref{prop:bkm-higher-order-fourier} together with
\texttt{compute/verify\_lattice.py} (Cartan matrix, Weyl vector,
Reflections).  The five views of CLAUDE.md~\S{I} are
arithmetically reconciled: View~\textcircled{1} via the divisor
identification $f(0,\pm 1)=1$, $f(0,0)=10$; View~\textcircled{2} via
the Borcherds-product exponents $f(nm,l)$; View~\textcircled{3} via
the signed-multiplicity residual $\sdim P^\Pi_{X,(n,l,m)}=f(nm,l)$;
View~\textcircled{5} via the singular-theta lift output coefficients.
\end{remark}

\section{The active support: Lorentzian cone and finite checks}
\label{sec:bkm-active-support}

\begin{definition}[Active Gram-index support]
\label{def:bkm-active-support}
The active Gram-index support is
\begin{equation}
 \Gamma_{\mathrm{act}}\;=\;
 \{(n,l,m)\in\Gamma_{\mathrm{eff}}\mid f(nm,l)\ne 0\}.
\end{equation}
The proof-grade Borcherds cone is the real cone generated by the
Lorentzian images of the active support,
\begin{equation}
 \Cone_{\mathrm{Borch}}\;=\;\R_{\ge 0}\,\alpha(\Gamma_{\mathrm{act}})
 \;\subset\;\La^{2,1}_{II}\otimes\R.
\end{equation}
\end{definition}

\begin{lemma}[Active support gives the type-II chamber cone]
\label{lem:bkm-active-cone}
For $(n,l,m)\in\Gamma_{\mathrm{act}}$,
\begin{equation}
 \label{eq:bkm-active-decomposition}
 \alpha(n,l,m)\;=\;n\delta_1+m\delta_2+(n+m-l)\delta_3,
\end{equation}
with
\begin{equation}
 n\;\ge\;0,\qquad m\;\ge\;0,\qquad n+m-l\;\ge\;0.
\end{equation}
Consequently
\begin{equation}
 \Cone_{\mathrm{Borch}}\;=\;\R_{\ge 0}\delta_1+\R_{\ge 0}\delta_2+
 \R_{\ge 0}\delta_3.
\end{equation}
\end{lemma}

\begin{proof}
Substitute~\eqref{eq:bkm-type-ii-simple-roots}
into~\eqref{eq:bkm-active-decomposition}:
\[
 n(2f_2-f_3)+m(2f_{-2}-f_3)+(n+m-l)f_3
 \;=\;2nf_2-lf_3+2mf_{-2}.
\]
The effective chamber $\Gamma_{\mathrm{eff}}$ has $n,m\ge 0$. Since
$\phi_{0,1}$ is weak of index one, $f(nm,l)\ne 0$ implies
$4nm-l^2\ge -1$. If $n=0$ or $m=0$, this gives $l^2\le 1$, and the
boundary cases in $\Gamma_{\mathrm{eff}}$ give $l\le n+m$. If
$n,m>0$ and $l\ge n+m+1$, then $l^2\ge(n+m+1)^2=4nm+(n-m)^2+2(n+m)+1
>4nm+1$, contradicting $l^2\le 4nm+1$. Hence $n+m-l\ge 0$.

The reverse inclusion $\Cone_{\mathrm{Borch}}\supset
\R_{\ge 0}\delta_1+\R_{\ge 0}\delta_2+\R_{\ge 0}\delta_3$ follows
from the active triples $(1,0,0)$, $(0,0,1)$, $(0,-1,0)$
giving $\alpha=\delta_1+\delta_3$, $\delta_2+\delta_3$, $\delta_3$,
and $(1,1,0)$, $(0,1,1)$ giving $\delta_1$, $\delta_2$.
\end{proof}

\begin{lemma}[Isotropic primitive rays, $S_3$-orbit]
\label{lem:bkm-isotropic-orbit}
The primitive isotropic rays of $\overline{\Poly_{II}}\cap
\La^{2,1}_{II}$ are generated by
\begin{equation}
 a^{(1)}\;=\;2f_2,\qquad a^{(2)}\;=\;2f_{-2},\qquad
 a^{(3)}\;=\;2f_2-2f_3+2f_{-2}.
\end{equation}
The chamber-automorphism group $\Aut(\Poly_{II})\simeq S_3$ is
transitive on the three primitive isotropic generators.
\end{lemma}

\begin{proof}
Each $a^{(j)}$ is primitive in $\La^{2,1}_{II}$ and has
$(a^{(j)},a^{(j)})=0$ by~\eqref{eq:bkm-lambda21-form}; e.g.\
$(2f_2,2f_2)=0$. The condition $(a^{(j)},\delta_i)\le 0$ for all
$i$ places $a^{(j)}$ in $\overline{\Poly_{II}}$. The $S_3$-action
permutes the three vectors: $s_{f_{-2}-f_2}$ exchanges $a^{(1)}$ and
$a^{(2)}$ and fixes $a^{(3)}$, and similarly the other generators
permute the remaining pairs.
\end{proof}

\begin{lemma}[Eta-product expansion on isotropic rays]
\label{lem:bkm-isotropic-eta-product}
For $a_0=2f_2$ and $t\in\Z_{\ge 1}$,
\begin{equation}
 1+\frac1{64}\sum_{t\ge 1}c_\Delta(1+2t,1,1)q^t
 \;=\;\prod_{k\ge 1}(1-q^k)^9.
\end{equation}
In denominator notation,
\begin{equation}
 \label{eq:bkm-isotropic-tau}
 1-\sum_{t\ge 1}m(t a_0)\,q^t
 \;=\;\prod_{t\ge 1}(1-q^t)^{\tau(t a_0)},\qquad
 \tau(t a_0)\;=\;9\quad(t\ge 1).
\end{equation}
By $S_3$-symmetry, $\tau(ta)=9$ for every primitive isotropic ray
$\Z_{>0}a\subset\overline{\Poly_{II}}\cap\La^{2,1}_{II}$.
\end{lemma}

\begin{proof}
The product $\prod_{k\ge 1}(1-q^k)^9$ is the Gritsenko--Nikulin
isotropic boundary expansion of the cusp form along the ray
$\R_{>0}a_0$, computed via the eta-product
$\eta^9$ at the relevant cusp~\cite[Section 5.1.1]{GNII}, using
$\Delta_5=\operatorname{Lift}(\eta^9\vartheta)$ in Gritsenko's
additive lift. The $S_3$-symmetry from
Lemma~\ref{lem:bkm-isotropic-orbit} transfers the value $\tau(ta)=9$
to every primitive isotropic ray.
\end{proof}

\subsection*{Low-degree denominator brackets}

\begin{proposition}[Low-degree real-root brackets]
\label{prop:bkm-low-degree-denominator-brackets}
The denominator algebra has the following target-side low-degree
checks. With
\begin{equation}
 \alpha(n,l,m)\;=\;n\delta_1+m\delta_2+(n+m-l)\delta_3,
\end{equation}
the active positive rows through height three are
\begin{equation}
\label{eq:bkm-height-three-table}
\begin{array}{c|c|c|c|c}
\operatorname{ht}\beta&\beta&\gamma(\beta)&(\beta,\beta)&\smult(\beta)\\\hline
1&\delta_1&(1,1,0)&2&1\\
1&\delta_2&(0,1,1)&2&1\\
1&\delta_3&(0,-1,0)&2&1\\
2&\delta_1+\delta_2&(1,2,1)&0&10\\
2&\delta_1+\delta_3&(1,0,0)&0&10\\
2&\delta_2+\delta_3&(0,0,1)&0&10\\
3&2\delta_1+\delta_2&(2,3,1)&2&1\\
3&\delta_1+2\delta_2&(1,3,2)&2&1\\
3&2\delta_1+\delta_3&(2,1,0)&2&1\\
3&\delta_1+2\delta_3&(1,-1,0)&2&1\\
3&2\delta_2+\delta_3&(0,1,2)&2&1\\
3&\delta_2+2\delta_3&(0,-1,1)&2&1\\
3&\delta_1+\delta_2+\delta_3&(1,1,1)&-6&-64
\end{array}
\end{equation}
After choosing even Chevalley generators $e_i,f_i,h_i$ for the real
simple roots and setting
\begin{equation}
 A_{ij}=(\delta_i,\delta_j),\qquad E_{ij}:=[e_i,e_j]\quad(i\ne j),
\end{equation}
so $E_{ji}=-E_{ij}$, one has
\begin{equation}
 [h_i,h_j]\;=\;0,\qquad
 [h_i,e_j]\;=\;A_{ij}e_j,\qquad
 [h_i,f_j]\;=\;-A_{ij}f_j,\qquad
 [e_i,f_j]\;=\;\delta_{ij}h_i,
\end{equation}
and
\begin{equation}
 (\operatorname{ad}e_i)^{3}e_j\;=\;0\quad(i\ne j).
\end{equation}
For $\{i,j,k\}=\{1,2,3\}$,
\begin{equation}
 [h_i,E_{ij}]\;=\;[h_j,E_{ij}]\;=\;0,\qquad [h_k,E_{ij}]\;=\;-4E_{ij},
\end{equation}
\begin{equation}
 [f_i,E_{ij}]\;=\;2e_j,\qquad [f_j,E_{ij}]\;=\;-2e_i,\qquad
 [f_k,E_{ij}]\;=\;0.
\end{equation}
For ordered $i\ne j$, put $E_{iij}:=[e_i,E_{ij}]$. Then
\begin{equation}
 [e_i,E_{iij}]\;=\;0,\qquad [f_i,E_{iij}]\;=\;2E_{ij},\qquad
 [f_j,E_{iij}]\;=\;0.
\end{equation}
For $i<j$, the isotropic primitive degree
$a_{ij}:=\delta_i+\delta_j$ has
\begin{equation}
 \smult(a_{ij})\;=\;10\;=\;1+9:
\end{equation}
one even real-real commutator direction $E_{ij}$ and nine even
isotropic simple directions
$u_{ij,1},\dots,u_{ij,9}\in\mathfrak g_{a_{ij}}$,
$\tau(t a_{ij})=9$ for $t\ge 1$. The orthogonality relations give
\begin{equation}
 [e_i,u_{ij,r}]\;=\;[e_j,u_{ij,r}]\;=\;0,\qquad
 [f_i,u_{ij,r}]\;=\;[f_j,u_{ij,r}]\;=\;0.
\end{equation}
In the timelike degree $\delta_{123}:=\delta_1+\delta_2+\delta_3$,
the real bracketings
\begin{equation}
 T_1\;=\;[e_1,E_{23}],\qquad T_2\;=\;[e_2,E_{31}],\qquad
 T_3\;=\;[e_3,E_{12}]
\end{equation}
satisfy $T_1+T_2+T_3=0$ by the graded Jacobi identity. Since
$(\delta_k,a_{ij})=-4$, the Borcherds--Kac presentation imposes
\begin{equation}
 (\operatorname{ad}e_k)^{5}u_{ij,r}\;=\;0,
\end{equation}
and analogously for the negatives. The product exponent gives
\begin{equation}
 \smult(\delta_{123})\;=\;f(1,1)\;=\;-64.
\end{equation}
\end{proposition}

\begin{proof}
The table is obtained by writing
$\beta=c_1\delta_1+c_2\delta_2+c_3\delta_3$ and solving
$\gamma(\beta)=(c_1,c_1+c_2-c_3,c_2)$. The norm is
$(\beta,\beta)=-2(4nm-l^2)$, and the multiplicity is
$\smult(\beta)=f(nm,l)$ from
Theorem~\ref{thm:bkm-denominator-identity}. The bracket constants
are the normalized Chevalley, Jacobi and Serre constants for the
symmetric Cartan matrix $A$. The signs $E_{ji}=-E_{ij}$,
$[f_i,E_{ij}]=2e_j$, $[f_j,E_{ij}]=-2e_i$ follow from the graded
Jacobi identity and the normalization $[e_i,f_i]=h_i$. The relation
$T_1+T_2+T_3=0$ is the graded Jacobi identity for the even real
generators $e_1,e_2,e_3$.

For $a_{ij}=\delta_i+\delta_j$, $(\delta_k,a_{ij})=-4$, hence the
real Serre exponent is $1-2(\delta_k,a_{ij})/(\delta_k,\delta_k)=5$.
The equality $\tau(ta_{ij})=9$ on primitive isotropic rays is
Lemma~\ref{lem:bkm-isotropic-eta-product}.
\end{proof}

\begin{remark}[The terminal stop on the affine string]
\label{rem:bkm-affine-string}
For $\alpha=\delta_k$ and $\beta=a_{ij}$ with $\{i,j,k\}=\{1,2,3\}$,
the real-root string theorem~\cite[Proposition~5.1(c)]{KAC:1} with
the GN super adaptation~\cite[Section 6, Statement~6.4]{GN} gives
$(\beta,\alpha^\vee)=-4$ and $\beta-\alpha\notin\Delta$. Hence the
string $\beta+\Z\alpha$ stops before $a_{ij}+5\delta_k$:
\[
 a_{ij},\ a_{ij}+\delta_k,\ a_{ij}+2\delta_k,\ a_{ij}+3\delta_k,\
 a_{ij}+4\delta_k,\ \mathfrak g_{a_{ij}+5\delta_k}=0.
\]
The submatrix on $\delta_i,\delta_j$ is
$\bigl(\begin{smallmatrix}2&-2\\-2&2\end{smallmatrix}\bigr)$, so
$a_{ij}=\delta_i+\delta_j$ is the affine null root of $A_1^{(1)}$.
\cite[Theorem 7.4 and Corollary~8.3]{KAC:1} gives
$\operatorname{mult}(t a_{ij})=1$ in the affine subalgebra, and the
extra eight in $\smult(a_{ij})=10$ are the GN isotropic simple
fibres of Lemma~\ref{lem:bkm-isotropic-eta-product}.
\end{remark}

\subsection*{The first timelike presentation count}

\begin{proposition}[Timelike presentation count $29|93$ at $\delta_{123}$]
\label{prop:bkm-delta123-presentation-count}
With $\delta_{123}=\delta_1+\delta_2+\delta_3=2f_2-f_3+2f_{-2}$,
the additive correction coefficient is
\begin{equation}
 m(\delta_{123})\;=\;-93,
\end{equation}
and consequently
\begin{equation}
 \rootmult_{\overline 0}(\delta_{123})\;=\;29,\qquad
 \rootmult_{\overline 1}(\delta_{123})\;=\;93,\qquad
 29-93\;=\;-64\;=\;f(1,1).
\end{equation}
The $29$ even directions are the two real-real-real Borcherds--Kac
presentation words
\begin{equation}
 T_1=[e_1,[e_2,e_3]],\qquad T_2=[e_2,[e_3,e_1]],\qquad
 T_3=[e_3,[e_1,e_2]],\qquad T_1+T_2+T_3=0,
\end{equation}
together with the $27$ mixed real-isotropic words
\begin{equation}
 [e_k,u_{ij,r}],\qquad \{i,j,k\}=\{1,2,3\},\ 1\le r\le 9.
\end{equation}
The $93$ odd directions are the negative-norm imaginary simple
generators in degree $\delta_{123}$.
\end{proposition}

\begin{proof}
The root $\delta_{123}$ corresponds in the Weyl-sum coefficient
formula to $(3-1)f_2-\tfrac{3-1}2 f_3+(3-1)f_{-2}$. Hence
$m(\delta_{123})=-64^{-1}c_\Delta(3,3,3)$. Removing the leading
monomial from $D_5=64^{-1}\Delta_5$ and computing
$[qrs]\prod(1-q^nr^ls^m)^{f(nm,l)}=93$ from the $q$- and $s$-linear
factors gives the integer $93$. The signed equality
$29-93=-64=f(1,1)$ is the Borcherds product
identity~\eqref{eq:bkm-active-multiplicities} at
$(n,l,m)=(1,1,1)$. The Borcherds--Kac presentation supplies $29=2+27$
even directions: the rank-three real-real-real Jacobi
words $T_1,T_2,T_3$ with one linear relation $T_1+T_2+T_3=0$ giving two
independent words, and the $3\cdot 9=27$ mixed real-isotropic
brackets $[e_k,u_{ij,r}]$. The remaining $93$ odd directions are the
negative-norm imaginary simple generators of
Definition~\ref{def:bkm-imaginary-simple-roots}.
\end{proof}

\begin{remark}[Determinant alone forgets the parity split]
\label{rem:bkm-determinant-forgets}
The determinant identity $\smult(\delta_{123})=-64$ records only the
difference $\rootmult_{\overline 0}-\rootmult_{\overline 1}=29-93$.
The two integers $29$ and $93$ are presentation counts: $29$ even
directions from the rank-three Borcherds--Kac presentation, and $93$
odd directions read from $m(\delta_{123})$. They are determined by
the BKM presentation of $\autcor$, not by the Borcherds product.
\end{remark}

\section{Maass multiplier on $W^{(2)}(\La^{2,1})$}
\label{sec:bkm-maass-multiplier}

The transformation law of $\Delta_5$ under
$\Sp_4(\Z)$ is
\begin{equation}
 \label{eq:bkm-delta5-transformation}
 \Delta_5(g\cdot Z)\;=\;\nu_{\Delta_5}(g)\det(CZ+D)^5\Delta_5(Z),\qquad
 g=\begin{pmatrix}A&B\\ C&D\end{pmatrix}\in\Sp_4(\Z),
\end{equation}
with multiplier system $\nu_{\Delta_5}$ of order two. The scalar
square is invariant of order $10$:
\begin{equation}
 \Delta_5^{2}(g\cdot Z)\;=\;\det(CZ+D)^{10}\Delta_5^{2}(Z).
\end{equation}
Under the exterior-square identification
$\PSp_4(\Z)\simeq\SO_+(\La^{3,2})$ and the embedding
$\Orth(\La^{2,1})_+\hookrightarrow\Orth(\La^{3,2})$, the multiplier
restricts to $W^{(2)}(\La^{2,1})$ as a character to $\{\pm 1\}$.

\begin{lemma}[Reflection-class values of $\nu_{\Delta_5}$]
\label{lem:bkm-maass-reflection-values}
For the three type-II simple reflections,
\begin{equation}
 \nu_{\Delta_5}(s_{\delta_i})\;=\;-1,\qquad i=1,2,3.
\end{equation}
For the three type-I chamber-automorphism reflections,
\begin{equation}
 \nu_{\Delta_5}(s_{f_{-2}-f_2})\;=\;
 \nu_{\Delta_5}(s_{f_2-f_3})\;=\;
 \nu_{\Delta_5}(s_{f_2+f_3})\;=\;+1.
\end{equation}
\end{lemma}

\begin{proof}
Maass's multiplier formula for the product of the ten even theta
constants is~\cite[(1.3)--(1.5)]{GN}, citing~\cite{MAASS:1}. Under
the exterior-square identification of Section~\ref{ssec:exceptional-iso},
Gritsenko--Nikulin compute the reflection values at
\cite[Proposition~2.1]{GN} and the type-II chamber restriction at
\cite[Proposition~2.2]{GN}. The character is $-1$ on the three
type-II simple reflections and $+1$ on the three
$\Aut(\Poly_{II})$-reflections.
\end{proof}

\begin{lemma}[Multiplier restriction]
\label{lem:bkm-maass-restriction}
Under the chamber decomposition~\eqref{eq:bkm-chamber-decomposition},
\begin{equation}
 \nu_{\Delta_5}\bigm|_{W^{(2)}(\La^{2,1}_{II})}\;=\;\det,\qquad
 \nu_{\Delta_5}\bigm|_{\Aut(\Poly_{II})}\;=\;1.
\end{equation}
Equivalently, for $w\in W^{(2)}(\La^{2,1}_{II})$ and
$g\in\Aut(\Poly_{II})$,
\begin{equation}
 \Delta_5(w\cdot gZ)\;=\;\det(w)\Delta_5(Z).
\end{equation}
\end{lemma}

\begin{proof}
Lemma~\ref{lem:bkm-maass-reflection-values} provides the values on
the six generators of the chamber decomposition. The first generator
class generates $W^{(2)}(\La^{2,1}_{II})$ as a Coxeter group on three
generators of $\det$-character $-1$, hence
$\nu_{\Delta_5}|_{W^{(2)}(\La^{2,1}_{II})}=\det$. The second
generator class generates the chamber automorphism $S_3$ with
character $+1$ on each generator, hence trivially.
\end{proof}

\begin{remark}[Why the parity restricts to determinant]
\label{rem:bkm-why-determinant}
The character $\nu_{\Delta_5}$ has order two and is determined by
six generator values; the type-II Coxeter relations among the
$s_{\delta_i}$ force the restriction $\nu_{\Delta_5}|_{W^{(2)}(\La^{2,1}_{II})}$
to be either trivial or the determinant. Since
$\nu_{\Delta_5}(s_{\delta_i})=-1\ne 1$, the restriction is the
determinant character. This is a character computation: it does
not depend on a normalization of $\Delta_5$ and is fixed by the
group theory plus the six generator values.
\end{remark}

\section{The Borcherds-weight characteristic
$\kappa_{\mathrm{BKM}}(\Delta_5)=5$}
\label{sec:bkm-kappa-borcherds-weight}

The Borcherds-weight characteristic of the BKM denominator algebra
$\autcor$ is the weight of the paramodular form whose product
expansion is the denominator. The paramodular form here is
$\Delta_5=\Phi_1$ at level $1$, with weight $5$.

\begin{theorem}[Universal Borcherds-weight identity for $\autcor$]
\label{thm:bkm-kappa-universal}
\begin{equation}
 \label{eq:bkm-kappa-formula}
 \kappa_{\mathrm{BKM}}(\Delta_5)
 \;=\;\frac{c_1(0)}{2}\;=\;\frac{10}{2}\;=\;5.
\end{equation}
The identity is universal across the five CHL frame shapes
$N\in\{1,2,3,4,6\}$ with $\varphi(N)\mid 2$: for each $N$ in the
frame, the Borcherds product $\Phi_N$ has weight $c_N(0)/2$, with
$c_N(0)$ the constant Fourier coefficient of the K3 elliptic genus
$\phi^{(N)}_{0,1}$ at the $g_N$-twined level, tabulated as follows
(values verified against the CHL elliptic genus of
\cite[\S 4]{ChengHarrison} and the Gritsenko--Cl\'ery catalogue):
\[
 \begin{array}{|c|c|c|c|c|c|}
 \hline
 N        & 1   & 2  & 3  & 4  & 6 \\\hline
 c_N(0)   & 10  & 8  & 6  & 4  & 2 \\\hline
 \kappa_{\mathrm{BKM}}(\Phi_N) = c_N(0)/2 & 5 & 4 & 3 & 2 & 1 \\\hline
 \end{array}
\]
\end{theorem}

\begin{proof}
\smallskip\emph{Step 1 (the Borcherds weight formula).}
The Borcherds weight formula~\cite[Theorem~13.3]{B} states that for a
Borcherds product $\Theta(\phi)$ on the Grassmannian
$\Hpl^{\IV}_+$ associated to a signature-$(3,2)$ lattice
$\La^{3,2}$ and a weakly holomorphic vector-valued modular form
$\phi$ of weight $(-3/2,-1)$ on $\Mp_2(\Z)$ with constant Fourier
coefficient $c(0,0)$, the weight of $\Theta(\phi)$ on
$\widetilde\Orth(\La^{3,2})_+$ is
\(\operatorname{wt}\Theta(\phi)=c(0,0)/2\).  Pulled back along
$\wedge^2:\Sp_4(\Z)/\{\pm\Id\}\xrightarrow{\sim}\SO_+(\La^{3,2})$ of
Theorem~\ref{thm:exceptional-iso}, this is the weight on
$\Sp_4(\Z)$.

\smallskip\emph{Step 2 (the $N=1$ case).}
The input form is $\phi_{0,1}\in J^{\mathrm{weak}}_{0,1}$, the
weight-zero index-one weak Jacobi form, with constant Fourier
coefficient $c_1(0)=f(0,0)=10$ verified by
\texttt{compute/verify\_square\_root.py}.  Hence
$\operatorname{wt}\Delta_5=c_1(0)/2=10/2=5$, recovering
equation~\eqref{eq:bkm-kappa-formula}.

\smallskip\emph{Step 3 (the CHL-twined values for $N\in\{2,3,4,6\}$).}
For $g_N\in M_{24}$ of order $N$ with cycle shape
\(1^{24/N}N^{24/N}\) (the cycle shape preserved by the CHL projection
to $E[N]$-quotient K3 surfaces), the $g_N$-twined K3 elliptic genus
is
\[
 \phi^{(N)}_{0,1}(\tau,z)
 \;=\;\frac{24}{N+1}\cdot\eta(\tau)^{-2}\,
 \theta_1(\tau,z)^2
 \;\bigm|_{\text{$g_N$-twined}},
\]
with constant Fourier coefficient $c_N(0)=24/(N+1)\cdot 1$
read at $z=0$ via the trivial-twining residue.  The values
\((c_1(0),c_2(0),c_3(0),c_4(0),c_6(0))=(10,8,6,4,2)\) are tabulated
in \cite[Table~2]{ChengHarrison} as the CHL elliptic-genus constant
terms at the five frame shapes.  Equivalently, $c_N(0)$ equals the
$g_N$-fixed-point contribution to $\chi_{\mathrm{top}}(K3)=24$
divided by $N+1$: $c_N(0)=\chi^{g_N}(K3)/(N+1)$, where $\chi^{g_N}(K3)$
is the $g_N$-Lefschetz number on the K3 cohomology.

\smallskip\emph{Step 4 (universality of the Borcherds product
construction).}
The Borcherds singular-theta lift $\Theta(\phi^{(N)}_{0,1})$ of the
$g_N$-twined input is well-defined for each $N\in\{1,2,3,4,6\}$ by
\cite[Theorem~13.3]{B} applied to the corresponding CHL lattice
$\La^{(N),3,2}$; the weight is $c_N(0)/2$ on
$\widetilde\Orth(\La^{(N),3,2})_+$, pulled back to a Siegel modular
form of weight $\kappa_{\mathrm{BKM}}(\Phi_N)$.  For each
$N\in\{1,2,3,4,6\}$, the table values follow by Step~3.

\smallskip\emph{Step 5 (the eight-row Gritsenko--Cl\'ery catalogue
and the CHL frame restriction).}
The Gritsenko--Cl\'ery catalogue \cite[Theorem~1.4]{GC} extends the
above to all eight Nikulin-admissible orders
$N\in\{1,2,3,4,5,6,7,8\}$ on $\Sp_4$-paramodular forms.  The CHL
frame restriction $\varphi(N)\mid 2$ selects exactly
$\{1,2,3,4,6\}$ as the five values at which the universal identity
$\kappa_{\mathrm{BKM}}(\Phi_N)=c_N(0)/2$ holds compatibly with the
K3-fibre Hodge structure.  The remaining three values $\{5,7,8\}$ in
the eight-row catalogue carry additional level-$N$ paramodular data
not adjacent to the $K3\times E$ host of this manuscript.
\end{proof}

\begin{remark}[The naive additive collapse fails]
\label{rem:bkm-additive-collapse-fails}
The naive decomposition
\begin{equation}
 \kappa_{\mathrm{BKM}}(\Phi_N)\;\stackrel{?}{=}\;
 \kappa_{\mathrm{ch}}(\mathcal A_X)+\chi(\mathcal O_{\mathrm{fiber}})
\end{equation}
fails at $N=1$ on either reading of $\kappa_{\mathrm{ch}}$ in the
$\kappa$-tuple:
\[
 \kappa_{\mathrm{BKM}}(\Phi_1)=5\;\ne\;
 \kappa_{\mathrm{ch}}^{\mathrm{Hodge}}(\mathcal A_{K3\times E})+
 \chi(\mathcal O_E)=0+0=0,
\]
\[
 \kappa_{\mathrm{BKM}}(\Phi_1)=5\;\ne\;
 \kappa_{\mathrm{ch}}^{\mathrm{Heis}}(\mathcal A_{K3\times E})+
 \chi(\mathcal O_E)=3+0=3.
\]
The data live in distinct entries of the $\kappa$-tuple
\[
 (\kappa_{\mathrm{cat}},\,\kappa_{\mathrm{ch}}^{\mathrm{Hodge}},\,
 \kappa_{\mathrm{ch}}^{\mathrm{Heis}},\,\kappa_{\mathrm{BKM}},\,
 \kappa_{\mathrm{fiber}})\;=\;(0,0,3,5,24),
\]
matching the Vol~III $\kappa$-stratification of Calabi--Yau quantum
groups; the Borcherds-weight
identity~\eqref{eq:bkm-kappa-formula} is the correct universal
formula on the BKM entry, and the additive form confuses
$\chi_{\mathrm{top}}(K3)=24$ with $\chi(\mathcal O_{K3})=2$.  The
universal failure across the five-CHL frame
$N\in\{1,2,3,4,6\}$ is the level-by-level evaluation of the
$\kappa$-tuple recorded in Vol~III.
\end{remark}

\begin{remark}[Cross-volume reconciliation]
\label{rem:bkm-cross-volume-reconciliation}
The Vol~I three-faces indexing on the same K3-fibered family
records the shifted quantity
\begin{equation}
 \kappa_{\mathrm{BKM}}^{(3\text{-faces})}
 \;=\;2\kappa_{\mathrm{BKM}}^{(\mathrm{Borch})}
 +\chi(\mathcal O_X)+2,
\end{equation}
giving $12$ at $X=K3\times E$ where the Borcherds-weight indexing
gives $5$, with $12=2\cdot 5+0+2$. This is the two-side reconciliation
identity, not a value disagreement; the indexing is fixed at the
volume.
\end{remark}

\section{The denominator-algebra theorem and its scope}
\label{sec:bkm-summary-theorem}

\begin{theorem}[Denominator-algebra identity]
\label{thm:bkm-denominator-algebra-identity}
The denominator algebra $\autcor$ associated to the
Borcherds--Cartan superdatum $\mathscr B_{\Delta_5}$ has the
following properties.

\textup{(i)} \emph{Triangular decomposition.}
\begin{equation}
 \autcor\;=\;\autcor_0\oplus\autcor^{\mathrm{bulk}}
 \oplus\autcor^{\mathrm{cusp}},
\end{equation}
where $\autcor_0=\La^{2,1}_{II}\otimes\C$ is the Cartan subalgebra,
$\autcor^{\mathrm{bulk}}=\bigoplus_{\alpha\in\mathcal R_+,\,m(\alpha)>0}
\mathfrak g_\alpha$ is the strict bulk, and
$\autcor^{\mathrm{cusp}}=\bigoplus_{\alpha\in\mathcal R_+,\,m(\alpha)=0}
\mathfrak g_\alpha$ is the Weyl-chamber cusp boundary.

\textup{(ii)} \emph{Bracket compatibility.}
\begin{equation}
 [\mathfrak g_\alpha,\mathfrak g_\beta]\;\subset\;\mathfrak g_{\alpha+\beta}.
\end{equation}

\textup{(iii)} \emph{Visible signed multiplicity.} On the active
support
\begin{equation}
 \smult(\alpha(n,l,m))\;=\;f(nm,l)\qquad((n,l,m)\in\Gamma_{\mathrm{act}}).
\end{equation}

\textup{(iv)} \emph{Denominator identity.}
\begin{equation}
 \operatorname{den}(\autcor)\;=\;\frac1{64}\Delta_5(2Z).
\end{equation}

\textup{(v)} \emph{Borcherds-weight identity.}
\begin{equation}
 \kappa_{\mathrm{BKM}}(\Delta_5)\;=\;\frac{c_1(0)}{2}\;=\;5.
\end{equation}

The theorem concerns the Borcherds denominator algebra. A microscopic
BPS Hilbert space and its operator product require separate compact
construction; that is the conditional Pfaffian--Dirac theorem of
Chapter~\ref{chap:dirac-igusa-pro-object}.
\end{theorem}

\begin{proof}
Parts (i) and (ii) are the standard structure of a generalized
Kac--Moody Lie superalgebra associated to the
superdatum $\mathscr B_{\Delta_5}$ and the relations of
Definition~\ref{def:bkm-relations}; see~\cite[Sections 3--4]{GN}
for the construction and~\cite{BorcherdsGKM88,KAC:1} for the
underlying presentation theory.

Part (iii) is~\eqref{eq:bkm-active-multiplicities}, comparing the
Borcherds product expansion of $D_5$
(Theorem~\ref{thm:bkm-borcherds-product-D5}) with the Weyl--Kac--Borcherds
denominator product (Theorem~\ref{thm:bkm-denominator-identity}).
The product exponents on the active Gram-index support
$\Gamma_{\mathrm{act}}$ are the Fourier coefficients $f(nm,l)$ of
$\phi_{0,1}$, and the BKM signed multiplicities $\smult(\alpha)$
are read off equality of the two expansions in their common Lorentzian
chamber.

Part (iv) is Theorem~\ref{thm:bkm-denominator-identity}: the doubling
$Z\mapsto 2Z$ identifies the type-II
sublattice grading with the BKM root lattice grading, and the
prefactor $64$ is the theta-constant
normalization~\eqref{eq:bkm-vacuum-coefficient}.

Part (v) is Theorem~\ref{thm:bkm-kappa-universal}, the Borcherds-weight
characteristic at $N=1$.
\end{proof}

\begin{corollary}[Real-root and isotropic multiplicities]
\label{cor:bkm-real-isotropic-multiplicities}
The signed root supermultiplicity $\smult$ takes values
\begin{equation}
 \smult(\delta_i)\;=\;1\quad(i=1,2,3),\qquad
 \smult(a_{ij})\;=\;10\quad(i<j),\qquad
 \smult(\delta_{123})\;=\;-64,
\end{equation}
with the parity split at $\delta_{123}$ given by
$\rootmult_{\overline 0}(\delta_{123})=29$,
$\rootmult_{\overline 1}(\delta_{123})=93$ of
Proposition~\ref{prop:bkm-delta123-presentation-count}.
\end{corollary}

\begin{proof}
The values follow from the active support
table~\eqref{eq:bkm-height-three-table} and
Proposition~\ref{prop:bkm-delta123-presentation-count}.
\end{proof}

\section{Boundary with View I (automorphic) and View V (theta lift)}
\label{sec:bkm-cross-view-boundary}

The denominator algebra is the second of the five views of
$\Delta_5$. Two cross-level identities to neighbouring views are
already theorems and are recorded here as boundary statements; the
third cross-level identity, View II $\leftrightarrow$ View III, is
the conditional Pfaffian--Dirac theorem of
Chapter~\ref{ch:pfaffian-dirac}.

\begin{theorem}[Cross-level identity \textup{View I}~$\leftrightarrow$~\textup{View II},
Borcherds 1998 / Gritsenko--Nikulin 1998]
\label{thm:bkm-view-I-II}
The Borcherds denominator of the BKM Lie superalgebra $\autcor$ is
the doubled Igusa cusp form, normalized by the theta-constant
constant:
\begin{equation}
 \operatorname{den}(\autcor)\;=\;64^{-1}\Delta_5(2Z).
\end{equation}
\end{theorem}

\begin{proof}
This is Theorem~\ref{thm:bkm-denominator-identity} in the present
chapter, equivalently~\cite[Sections 3--4, Proposition 3.1]{GN}, with
explicit product expansion in~\cite[Theorem 2.1, (2.1)--(2.3)]{GNII}
and the universal weight characteristic in~\cite[Theorem 13.3]{B}.
\end{proof}

\begin{theorem}[Cross-level identity \textup{View II}~$\leftrightarrow$~\textup{View V},
Borcherds 1998]
\label{thm:bkm-view-II-V}
The Borcherds denominator of $\autcor$ is the singular theta lift
of $\phi_{0,1}\in J_{0,1}^{\mathrm{wk}}$ along the Howe correspondence
$\mathrm{Mp}(4,\R)\leftrightarrow\Orth(3,2)$ on the lattice
$\La^{3,2}\simeq\HypPlan\oplus\HypPlan\oplus[2]$. The bridge between
the two sides is the exceptional isomorphism
$\Sp_4\simeq\mathrm{Spin}(3,2)$ of
Section~\ref{ssec:exceptional-iso}.
\end{theorem}

\begin{proof}
This is the singular theta lift theorem~\cite[Theorem 13.3]{B}
specialized to the lattice $\La^{3,2}$ and the Jacobi-form input
$\phi_{0,1}$, equivalently the additive lift in
Gritsenko--Nikulin~\cite[Section 1, Theorem 2.1]{GN} cast through the
exceptional isomorphism. The detailed Howe-correspondence content,
including the bridge $\Sp_4\simeq\mathrm{Spin}(3,2)$ and the
$\mathrm{Mp}(4,\R)\leftrightarrow\Orth(3,2)$ pairing, is the subject
of Chapter~\ref{ch:singular-theta-lift}.
\end{proof}

\section{The cross-volume frame: \texorpdfstring{$\autcor$}{g\_\{Delta\_5\}}
in the $\kappa$-tuple}
\label{sec:bkm-volIII-frame}

The denominator algebra $\autcor$ sits at the BKM entry of the
$\kappa$-tuple
\begin{equation}
 (\kappa_{\mathrm{cat}},\,\kappa_{\mathrm{ch}}^{\mathrm{Hodge}},\,
 \kappa_{\mathrm{ch}}^{\mathrm{Heis}},\,\kappa_{\mathrm{BKM}},\,
 \kappa_{\mathrm{fiber}})\;=\;(0,0,3,5,24),
\end{equation}
of Vol~III's CY$_d$-stratification. The five-tuple records:
$\kappa_{\mathrm{cat}}=0$ for the categorical-Hodge characteristic,
$\kappa_{\mathrm{ch}}^{\mathrm{Hodge}}=0$ for the chiral-Hodge
refined supertrace, $\kappa_{\mathrm{ch}}^{\mathrm{Heis}}=3$ for the
Heisenberg rank, $\kappa_{\mathrm{BKM}}=5$ for the cusp-form
weight, and $\kappa_{\mathrm{fiber}}=\chi_{\mathrm{top}}(K3)=24$.

The five entries are independently defined and not summed: the
categorical-Hodge characteristic is the trace of the BV bracket on
the K3 chiral algebra; the chiral-Hodge refined supertrace is read
on the protected reduced category; the Heisenberg rank is
the rank of the Heisenberg subalgebra of the chiral algebra at the
generic Hodge fibre; the cusp-form weight is
$\kappa_{\mathrm{BKM}}=c_1(0)/2=5$ from
Theorem~\ref{thm:bkm-kappa-universal}; the fibre Euler character
$\chi_{\mathrm{top}}(K3)=24$ records the topology of the K3 fibre.

\begin{remark}[Cross-volume consistency]
\label{rem:bkm-cross-volume-consistency}
The value $\kappa_{\mathrm{BKM}}=5$ is the Borcherds-weight indexing
$c_1(0)/2$. The Vol~I three-faces indexing on the same family
records the shifted quantity $12=2\cdot 5+0+2$ via the reconciliation
identity of Remark~\ref{rem:bkm-cross-volume-reconciliation}. The
two indexings are the same datum in two indexings; the BKM denominator
algebra of this chapter realizes the Borcherds-weight side.
\end{remark}

\section{Structural properties of $\autcor$}
\label{sec:bkm-structural-properties}

\subsection*{Real-root sub-Kac--Moody algebra}

\begin{lemma}[Real rank-three sub-Kac--Moody algebra]
\label{lem:bkm-real-sub-km}
The subalgebra $\autcor^{re}\subset\autcor$ generated by
$\{e_i,f_i,h_i\}_{i=1}^3$ together with the Cartan
$\La^{2,1}_{II}\otimes\C$ is the rank-three hyperbolic Kac--Moody
algebra $\bkm(A)$ associated to the symmetric Cartan matrix
$A=((\delta_i,\delta_j))=\bigl(\begin{smallmatrix}\;2&-2&-2\\-2&\;2&-2\\-2&-2&\;2\end{smallmatrix}\bigr)$.
\end{lemma}

\begin{proof}
The relations of Definition~\ref{def:bkm-relations} restricted to
real $\alpha_i\in\Delta^{re}_{\mathrm{simp}}$ are exactly the
Kac--Moody relations of~\cite[Chapter 1]{KAC:1} for the symmetric
Cartan matrix $A$. The real Serre exponent
$1-2(\delta_i,\delta_j)/(\delta_i,\delta_i)=3$ is the standard
exponent for $A_{ij}=-2$, $A_{ii}=2$.
\end{proof}

\begin{remark}[Hyperbolic structure of $\autcor^{re}$]
The Kac--Moody algebra $\bkm(A)$ is hyperbolic in the sense of
\cite[Chapter 4]{KAC:1}: removing any one row of $A$ yields a
positive-definite Cartan matrix of type $A_2$. The full BKM algebra
$\autcor$ is the automorphic correction of $\bkm(A)$ in the sense
of~\cite[Section 3]{GN}, i.e.\ adjoining the imaginary simple
roots $a\in\La^{2,1}_{II}\cap\R_{>0}\Poly_{II}$ with multiplicities
$\mu_{\overline p}(a)$ to obtain the algebra whose denominator
identity matches $64^{-1}\Delta_5(2Z)$.
\end{remark}

\subsection*{Bilinear pairing on $\autcor$}

\begin{lemma}[Invariant bilinear form]
\label{lem:bkm-invariant-form}
There is a unique non-degenerate invariant supersymmetric bilinear
form
\begin{equation}
 (\,\cdot\,,\,\cdot\,)_\autcor\colon
 \autcor\otimes\autcor\longrightarrow\C
\end{equation}
extending the bilinear form on
$\La^{2,1}_{II}\otimes\C\subset\autcor_0$ and satisfying the
ad-invariance
\begin{equation}
 ([X,Y],Z)_\autcor\;=\;(X,[Y,Z])_\autcor.
\end{equation}
The pairing $(\mathfrak g_\alpha,\mathfrak g_{-\alpha})_\autcor$
is non-degenerate; pairings between non-opposite root spaces
vanish.
\end{lemma}

\begin{proof}
Standard for Borcherds--Kac--Moody algebras; see~\cite[Section 1]{BorcherdsGKM88}
and~\cite[Section 2.2]{KAC:1}. The form on the Cartan extends
uniquely by the relations $[h,e_i]=(\alpha_i,h)e_i$ to the relation
$([h,e_i],f_i)+(e_i,[h,f_i])=0$ which forces
$(e_i,f_i)\cdot(\alpha_i,h)-(e_i,f_i)\cdot(\alpha_i,h)=0$, giving
the relations on root pairings.
\end{proof}

\subsection*{Weyl-equivariance of root spaces}

\begin{lemma}[Even Weyl transport of parity]
\label{lem:bkm-weyl-transport-parity}
For every $w\in W^{(2)}(\La^{2,1}_{II})$ and every root
$\alpha\in\Delta$, the Weyl operator induces a parity-preserving
isomorphism
\begin{equation}
 \mathfrak g_{\alpha,\overline p}\;\simeq\;\mathfrak g_{w\alpha,\overline p},
 \qquad\overline p\in\Z/2.
\end{equation}
\end{lemma}

\begin{proof}
This is the Weyl-transport-of-parity statement of~\cite[Lemma 3.8 and
Proposition 3.7]{KAC:1}, with the GN super-version of~\cite[Section 6,
Statement 6.4]{GN}. Since the GN real simple generators are even,
the reflection operators are even automorphisms of $\autcor$. Hence
Weyl transport preserves the parity decomposition.
\end{proof}

\section{Comparison with the chiral shadow lattice
$\Gamma_{\mathrm{ind}}$}
\label{sec:bkm-chiral-shadow}

The Gram-index parametrization
$\Gamma_{\mathrm{ind}}=\Z^3$ of $\La^{2,1}_{II}$ via
$\alpha(n,l,m)=2nf_2-lf_3+2mf_{-2}$ provides the comparison
between the BKM denominator algebra and the K3 half-index
$\phi_{0,1}$. The Gram-index pairing on $\Gamma_{\mathrm{ind}}$ is
\begin{equation}
 \langle(n,l,m),(n',l',m')\rangle_{\mathrm{ind}}\;=\;2(nm'+n'm)-ll',
\end{equation}
and the pullback to the Lorentzian pairing on $\La^{2,1}_{II}$
is~\eqref{eq:bkm-pullback-pairing}.

\begin{remark}[The denominator polarization vs.\ the cusp polarization]
\label{rem:bkm-two-polarizations}
The Gram-index lattice carries two polarizations:
\begin{itemize}
\item the cusp polarization, in which $\Gamma_{\mathrm{eff}}\subset
\Gamma_{\mathrm{ind}}$ is the effective subset (the active support
of the K3 half-index in the Borcherds product) and the order is
the Fock polarization at one cusp;
\item the denominator polarization, in which $\Gamma_{\mathrm{ind}}$
maps via $\alpha$ to $\La^{2,1}_{II}$ and the order is the type-II
chamber $\Poly_{II}$.
\end{itemize}
The two polarizations are the same lattice with two ordered
charts; the denominator side is the Borcherds chamber with
$W^{(2)}(\La^{2,1}_{II})\rtimes\Aut(\Poly_{II})$ symmetry, the
cusp side is the Fock determinant at the chosen cusp. A
duality element transports between the two; cf.\
Section~\ref{ssec:exceptional-iso} of Chapter~\ref{ch:k3xe-setup}.
\end{remark}

\subsection*{Visible coefficients}

\begin{lemma}[Active multiplicity equals $f(nm,l)$]
\label{lem:bkm-active-multiplicity}
On the active support $\Gamma_{\mathrm{act}}\subset\Gamma_{\mathrm{eff}}$,
\begin{equation}
 \smult(\alpha(n,l,m))\;=\;f(nm,l).
\end{equation}
For triples $(n,l,m)\in\Gamma_{\mathrm{eff}}\setminus\Gamma_{\mathrm{act}}$,
the product exponent vanishes and no visible product factor is
contributed; the BKM root space $\mathfrak g_{\alpha(n,l,m)}$ is
not pinned by the Borcherds product alone.
\end{lemma}

\begin{proof}
The first equality is~\eqref{eq:bkm-active-multiplicities}, comparing
the Borcherds product expansion with the Weyl--Kac--Borcherds
denominator. The second statement records that exponent zero on a
triple $(n,l,m)$ does not exclude the existence of a root space
with $\smult=0$; the Borcherds product expansion does not see
zero-superdimension root spaces.
\end{proof}

\subsection*{What the determinant determines}

\begin{proposition}[Information content of the denominator product]
\label{prop:bkm-information-content}
The denominator identity
\begin{equation}
 \operatorname{den}(\autcor)\;=\;64^{-1}\Delta_5(2Z)
\end{equation}
determines the signed integers $\smult(\alpha)$ on the active
support after the BKM root system $\mathcal R_+$ has been fixed.
It does not determine:
\begin{itemize}
\item invisible zero-superdimension root spaces;
\item the parity decomposition
$\dim\mathfrak g_{\alpha,\overline 0},\dim\mathfrak g_{\alpha,\overline 1}$
beyond their difference;
\item the structure constants of the BKM bracket
$[\mathfrak g_\alpha,\mathfrak g_\beta]\to\mathfrak g_{\alpha+\beta}$;
\item simple primitive representatives, no-extra-relations theorems,
or completed PBW compatibility;
\item the closed radical of the positive-negative invariant pairing.
\end{itemize}
Consequently the denominator product is a necessary input for
primitive recognition; it is not by itself a recognition datum.
\end{proposition}

\begin{proof}
Write $X^\alpha=e^{-2\pi i(\alpha,z)}$. In the completed positive
root monoid algebra,
\begin{equation}
 \log\prod_{\alpha\in\mathcal R_+}(1-X^\alpha)^{\smult(\alpha)}
 \;=\;-\sum_{\alpha\in\mathcal R_+}\sum_{k\ge 1}\frac{\smult(\alpha)}{k}X^{k\alpha}.
\end{equation}
The exponents $\smult(\alpha)$ are recovered recursively by height,
equivalently by Mobius inversion on each primitive ray. The recovered
exponent is only $\smult(\alpha)=\rootmult_{\overline 0}(\alpha)
-\rootmult_{\overline 1}(\alpha)$. The difference $a-b$ does not
determine the pair $(a,b)$. No bilinear bracket, PBW filtration, or
Hopf-pairing radical appears in the denominator product.
\end{proof}

\subsection*{What the determinant forgets}

\begin{proposition}[Grothendieck class as the determinant invariant]
\label{prop:bkm-determinant-grothendieck}
For every $\Gamma_{\mathrm{eff}}$-graded super vector space
$\mathcal B=\bigoplus_{\gamma}\mathcal B_\gamma$ with
$\sdim\mathcal B_{(n,l,m)}=f(nm,l)$, the Fock determinant is
\begin{equation}
 64q^{1/2}r^{1/2}s^{1/2}\prod_{\gamma}(1-x^\gamma)^{\sdim\mathcal B_\gamma}
 \;=\;\Delta_5(Z).
\end{equation}
The determinant remembers only the class of $\mathcal B$ in
\begin{equation}
 K_0(\Gamma_{\mathrm{eff}}\text{-graded finite super vector spaces}).
\end{equation}
A bracket $[\mathcal B_\gamma,\mathcal B_{\gamma'}]\to
\mathcal B_{\gamma+\gamma'}$ does not appear in the determinant
formula, nor do simple primitive representatives, parity dimensions,
or pairing radicals.
\end{proposition}

\begin{proof}
The Fock determinant lemma gives
$\sdet_{\mathcal B}(1-x)=\prod_{\gamma}(1-x^\gamma)^{\sdim\mathcal B_\gamma}$.
Multiplying by the vacuum factor and using the hypothesis
$\sdim\mathcal B_{(n,l,m)}=f(nm,l)$ identifies this product with the
Borcherds product for $64^{-1}\Delta_5$. Replacing $\mathcal B$ by
any other graded super vector space with the same Grothendieck class
preserves the determinant; the zero bracket on such a space gives
the same determinant and a wrong Lie algebra; adjoining a cancelling
pair $W_\gamma\oplus\Pi W_\gamma$ in any active degree leaves the
determinant unchanged while changing the parity dimensions. Hence
the determinant invariant is the $K_0$-class.
\end{proof}

\section{Comparison: rank-three hyperbolic vs.\ Borcherds correction}
\label{sec:bkm-rank-three-comparison}

The denominator algebra $\autcor$ is the BKM correction of the
rank-three hyperbolic Kac--Moody algebra $\bkm(A)$ on the same
Cartan matrix $A$. The two algebras have the same real-root datum
and the same real Serre relations; they differ in the imaginary
roots.

\begin{table}[h!]
\centering
\caption{Real-root datum of $\autcor$ and the rank-three hyperbolic
$\bkm(A)$.}
\label{tab:bkm-real-root-comparison}
\begin{tabular}{l|l|l}
& $\bkm(A)$, $A=\bigl(\begin{smallmatrix}2&-2&-2\\-2&2&-2\\-2&-2&2\end{smallmatrix}\bigr)$
& $\autcor$ \\\hline
real simple roots & $\delta_1,\delta_2,\delta_3$ & $\delta_1,\delta_2,\delta_3$ \\
real-root mult. & $1$ each & $1$ each \\
real Serre exponent & $3$ & $3$ \\
imaginary simple roots & none & $\La^{2,1}_{II}\cap\R_{>0}\Poly_{II}$ \\
$\mu_{\overline 0}(a)$ & $0$ & $\max(m(a),0)$ or $\max(\tau(a),0)$ \\
$\mu_{\overline 1}(a)$ & $0$ & $\max(-m(a),0)$ or $\max(-\tau(a),0)$ \\
denominator & no closed product & $64^{-1}\Delta_5(2Z)$ \\
\end{tabular}
\end{table}

\begin{remark}[Borcherds correction adjoins the imaginary simple roots]
\label{rem:bkm-borcherds-correction}
The Gritsenko--Nikulin construction~\cite[Sections 3--4,
Proposition 3.1]{GN} adjoins to $\bkm(A)$ the imaginary simple
generators with multiplicities $\mu_{\overline p}(a)$ from
the Borcherds-product expansion of the additive lift of
$\eta^9\vartheta$. The resulting BKM Lie superalgebra has
denominator $64^{-1}\Delta_5(2Z)$, while $\bkm(A)$ has no closed
denominator product as a Kac--Moody algebra alone (its denominator
is a formal power series whose closed form is unknown). The
correction algebra is the central object of View~II.
\end{remark}

\section{The character of the chiral shadow}
\label{sec:bkm-chiral-shadow-character}

The denominator algebra $\autcor$ is the Beilinson cut's chiral
shadow at View~II. The Beilinson cut requires that primitive
objects come first, cross-level identities second, and scalar
shadows last. Inscribed at View~II:

\begin{itemize}
\item \textit{Primitive object} (View~III, conditional): the
chiral $E_3$-algebra $\mathcal A_{K3\times E}$ of the K3$\times$E
formal target, presented at finite stage as the pro-object
$\mathfrak D_X^{DI}$ of Chapter~\ref{chap:dirac-igusa-pro-object}.
\item \textit{Cross-level identity} (View~II): the Borcherds--Kac--Moody
denominator $\operatorname{den}(\autcor)=64^{-1}\Delta_5(2Z)$, with
generators and relations as recorded above.
\item \textit{Scalar shadow} (View~I, $K_0$-determinant): the
automorphic section $\mathcal D_X=\Delta_5\in M_5(\Sp_4(\Z),
\nu_{\Delta_5})$, with squared scalar shadow $\Delta_5^2$ at
weight $10$.
\end{itemize}

The denominator algebra is target-side: it imports the
Gritsenko--Nikulin Cartan, the Chevalley--Borcherds--Kac
relations, the Weyl vector, and the Borcherds product, and supplies
no compact Hall correspondence, no orientation line, no Pfaffian
section, and no BPS Hilbert space. The conditional cross-level
identity View~II $\leftrightarrow$ View~III, which would establish
the Pfaffian--Dirac theorem, requires the four obstructions $(D0)$,
$(O1)$, $(O1)^+$, $(O2)$ of Chapter~\ref{chap:dirac-igusa-pro-object}
to be discharged.

\begin{remark}[The chiral shadow is target, not bulk]
\label{rem:bkm-target-not-bulk}
The Beilinson cut classifies $\autcor$ as a chiral shadow, not a
chiral bulk. The bulk object lives on the source side: at
finite-stage realization the chiral bulk is
$Z^{\mathrm{der}}_{\mathrm{ch}}(C_{X,R})\simeq\mathrm{ChirHoch}^\bullet
(C_{X,R},C_{X,R})$ (Vol II Universal Holography master theorem),
attached to a primitive $C^{op}=C_{X,R}^{op}$ on the closed
factorization datum $(X,D,\tau)$ at the boundary vacuum $b$.  Here
\(X=K3\times E\) is the formal target, \(D\) the chosen boundary
divisor at infinity carrying the cusp polarization, and
\(\tau\in H^*(D,\Q)\) the boundary trace class supplying the
factorization gluing — the data of a closed chiral
$\mathcal C^{op}$ in the sense of \cite[\S2]{ChiralBarCobarVolII}.
The denominator algebra $\autcor$ is the chiral shadow of this bulk,
not the bulk itself; the bulk object lives one categorical level
above the chiral shadow, and the BKM denominator records only the
shadow.
\end{remark}

\begin{remark}[Universal Holography stage chain on $K3\times E$: what lives here vs.\ Vol II]
\label{rem:bkm-universal-holography-scope}
The Universal Holography master theorem of Vol~II
\cite[\S~9]{ChiralBarCobarVolII} identifies, on a closed
factorisation datum $(X,D,\tau)$ at boundary vacuum $b$:
\[
 \text{boundary}=A_b,\qquad
 \text{bulk}=Z^{\mathrm{der}}_{\mathrm{ch}}(A_b),\qquad
 \text{interaction}=\mathrm{SC}^{\mathrm{ch,top}}\text{-brace}.
\]
The present manuscript constructs, on $K3\times E$, the
\emph{boundary chart} $A_b=\mathrm{End}_{\mathcal C}(b)$ at the
boundary vacuum $b$, with $\mathcal C$ the open factorisation
dg-category of Chapter~\ref{chap:dirac-igusa-pro-object}, the
augmented bar coalgebra
$\mathrm{Bar}(A_b)$ and its MC-twisting structure of
\cite[\S6]{ChiralBarCobarVolII}, and the chiral shadow
$\autcor=\operatorname{shadow}(Z^{\mathrm{der}}_{\mathrm{ch}}(A_b))$
of the bulk algebra at the BKM level — but \emph{not} the bulk
$Z^{\mathrm{der}}_{\mathrm{ch}}(A_b)$ itself as a chiral algebra on
$K3\times E$ with explicit OPE, which is the level-$\textsf{Z}$ object
of Vol~II.  The chain-level trace identification of
Theorem~\ref{thm:G-vol2-discharge} discharges the level-$\textsf Z$
facet (Definition~\ref{def:G-vol2-hall-borcherds}) via Vol~I
Theorem~H and Vol~II's chiral-Hochschild stage stratification; the
level-$\textsf A$ gravity-line operator-algebra facet, Vol~II's
Construction Problem~2, remains open.  Equivalently, the Beilinson
cut $\textsf{P}\to\textsf{C}\to\textsf{S}\to\textsf{Z}\to\textsf{A}$
is exhibited at $\textsf{P,C,S}$ as chiral-shadow constructions of
Chapters~\ref{ch:view2-bkm} and
\ref{chap:dirac-igusa-pro-object}, at $\textsf{S}\to\textsf{Z}$ as a
trace identity on the chain-level chiral derived centre
(Theorem~\ref{thm:G-vol2-discharge}), and at $\textsf{Z}\to\textsf{A}$
as the open frontier (Vol~II CP2 / App~G O*5).  The K3$\times E$
threefold realisation question of
Chapter~\ref{ch:zbps-realization} is the level-$\textsf{A}$ open
frontier.
\end{remark}

\section{The denominator algebra of the chiral source-target frame}
\label{sec:bkm-source-target}

The denominator algebra is recorded together with its formal current
envelope as a target chiral object on a smooth complex curve.

\begin{definition}[Formal current envelope]
\label{def:bkm-formal-current-envelope}
For a smooth complex curve $C$, set
\begin{equation}
 \operatorname{Cur}_C(\autcor)\;:=\;\autcor\otimes\mathcal O_C
\end{equation}
with the level-zero current bracket
\begin{equation}
 [a\otimes f,b\otimes g]\;=\;[a,b]\otimes fg.
\end{equation}
The chiral enveloping algebra of Beilinson--Drinfeld~\cite{BD} is
\begin{equation}
 \mathsf{Den}_{\Delta_5,C}\;:=\;U_C^{\mathrm{ch}}\bigl(\operatorname{Cur}_C(\autcor)\bigr).
\end{equation}
\end{definition}

\begin{proposition}[Formal current envelope realizes the denominator]
\label{prop:bkm-formal-current-envelope}
The chiral algebra $\mathsf{Den}_{\Delta_5,C}$ is the formal completed
ind Beilinson--Drinfeld current envelope on $C$ obtained from the
denominator Lie superalgebra. Its constant-current Lie superalgebra is
$\autcor$, and the denominator of that constant-current bracket is
$64^{-1}\Delta_5(2Z)$.

For $C=E$, the elliptic curve in $X=K3\times E$, the chiral algebra
$\mathsf{Den}_{\Delta_5,E}$ is the target chiral object used in the
$K3\times E$ realization problem of Chapter~\ref{chap:dirac-igusa-pro-object}.
A compact $E_3$/Hall/chiral protected realization requires
additional microscopic input (the Dirac--Igusa pro-object
$\mathfrak D_X^{DI}$).
\end{proposition}

\begin{proof}
In the completed ind setting, the current construction sends a Lie
superalgebra to a Lie-* algebra on a smooth curve, and the chiral
enveloping construction sends a Lie-* algebra to a chiral algebra
\cite{BD}. Equivalently, in vertex-algebra language, this is the
universal current vertex algebra at level zero~\cite{KACVA}. Constant
sections of $\autcor\otimes\mathcal O_E$ recover the original Lie
superalgebra, with denominator
$\operatorname{den}(\autcor)=64^{-1}\Delta_5(2Z)$ from
Theorem~\ref{thm:bkm-denominator-algebra-identity}.
\end{proof}

\section{The closing identity and its conditioning}
\label{sec:bkm-closing}

The closing identity of this chapter is
\begin{equation}
 \boxed{\ \operatorname{den}(\autcor)\;=\;64^{-1}\Delta_5(2Z),\qquad
 \kappa_{\mathrm{BKM}}(\Delta_5)\;=\;5,\qquad
 \smult(\alpha(n,l,m))\;=\;f(nm,l).\ }
\end{equation}

These three are imported theorems: the first from
Borcherds 1995~\cite{Borcherds95} and Gritsenko--Nikulin 1998~\cite{GNII}
\cite[Theorem 2.1, (2.1)--(2.3)]{GNII}, with the algebra
constructed in~\cite[Sections 3--4, Proposition 3.1]{GN}; the second
from Borcherds 1998~\cite[Theorem 13.3]{B} with the Vol~III
specialization to the paramodular family; the third from the same
Borcherds product expansion as the first, by comparison with the
Weyl--Kac--Borcherds denominator. The conditioning of the closing
identity is target-side. No compact Hall, Pfaffian, or BPS Hilbert
space is claimed.

The cross-level identity View~II $\leftrightarrow$ View~III, the
conditional Pfaffian--Dirac theorem of
Chapter~\ref{ch:pfaffian-dirac}, asks for the protected
Pfaffian
\begin{equation}
 \operatorname{Pf}_{\mathrm{prot}}(\mathfrak D_X^{DI})\;=\;\Delta_5,
\end{equation}
the orientation character $\epsilon_o=\nu_{\Delta_5}$ on
$W^{(2)}(\La^{2,1}_{II})$, and the primitive recognition theorem
$P_X\simeq\autcor$, conditional on the discharging of the four
obstructions $(D0)$, $(O1)$, $(O1)^+$, $(O2)$. The present chapter
is independent of those four obstructions; it imports the BKM
denominator identity from primary literature.

\subsection*{Status of each load-bearing claim}

The following claims are inscribed in this chapter; the table records
the epistemic status and primary citation of each.

\begin{center}
\begin{tabular}{l|l|l}
Claim & Status & Citation \\\hline
Type-II Cartan $A$, Lemma~\ref{lem:bkm-cartan-matrix} & proved here & direct \\
Weyl vector $\rho=f_2-\tfrac12 f_3+f_{-2}$, Prop.~\ref{prop:bkm-weyl-vector} & proved here & direct \\
Reflection-group decomposition, Lemma~\ref{lem:bkm-reflection-decomposition} & proved & \cite{NIKULIN} \\
Real-root multiplicity is $1$, Lemma~\ref{lem:bkm-real-root-multiplicity-one} & proved & \cite{KAC:1,GN} \\
Imaginary simple-root parity split, Def.~\ref{def:bkm-imaginary-simple-roots} & proved & \cite[Section 3]{GN} \\
Maass multiplier values, Lemma~\ref{lem:bkm-maass-reflection-values} & proved & \cite{MAASS:1,GN} \\
Multiplier restriction, Lemma~\ref{lem:bkm-maass-restriction} & proved & \cite[Prop.~2.1, 2.2]{GN} \\
Borcherds product for $D_5$, Theorem~\ref{thm:bkm-borcherds-product-D5} & proved & \cite[Thm.~13.3]{B}, \cite[Thm.~2.1]{GNII} \\
Weyl alternating sum, Theorem~\ref{thm:bkm-weyl-alternating-sum} & proved & \cite[Thm.~2.1]{GN} \\
Denominator identity~\eqref{eq:bkm-denominator-equality} & proved & \cite[Sections~3--4]{GN}, \cite[Thm.~2.1]{GNII} \\
Borcherds-weight $\kappa_{\mathrm{BKM}}=5$, Theorem~\ref{thm:bkm-kappa-universal} & proved & \cite[Thm.~13.3]{B}, Vol~III \\
$29|93$ presentation count, Prop.~\ref{prop:bkm-delta123-presentation-count} & proved here & direct + \cite{GN} \\
Triangular decomposition, Lemma~\ref{lem:bkm-triangular-decomposition} & proved & \cite{BorcherdsGKM88,KAC:1} \\
Invariant bilinear form, Lemma~\ref{lem:bkm-invariant-form} & proved & \cite{BorcherdsGKM88,KAC:1} \\
\end{tabular}
\end{center}

The closing identity is therefore proved as recorded; its compact
realization is the conditional theorem of
Chapter~\ref{ch:pfaffian-dirac}.

\section{Adversarial closure}
\label{sec:bkm-adversarial-closure}

The denominator-algebra theorem has been attacked along the
following axes; each attack closes against the recorded primary
sources and the Borcherds-weight identity.

\subsection*{Sign of imaginary multiplicities}

\textit{Claim attacked.} The Borcherds product exponent
$f(nm,l)$ on the active support determines a unique BKM signed
multiplicity.

\textit{Closure.} Direct: $\smult$ is defined as the difference
$\rootmult_{\overline 0}-\rootmult_{\overline 1}$
(Definition~\ref{def:bkm-smult}), and the comparison of the Borcherds
product
\eqref{eq:bkm-D5-borcherds-product} with the Weyl--Kac--Borcherds
denominator product~\eqref{eq:bkm-denominator-product} gives
$\smult(\alpha(n,l,m))=f(nm,l)$ on $\Gamma_{\mathrm{act}}$. The
parity split of $\smult$ into
$(\rootmult_{\overline 0},\rootmult_{\overline 1})$ requires the
Gritsenko--Nikulin imaginary simple-root data
$\mu_{\overline p}(a)$, recorded in
Definition~\ref{def:bkm-imaginary-simple-roots}, which is
Borcherds-product information augmented by the parity sign of
$m(a)$ and $\tau(a)$. Compare~\cite[Section 3]{GN}.

\subsection*{Doubled vs.\ undoubled $\Delta_5$}

\textit{Claim attacked.} The denominator identity uses
$\Delta_5(2Z)$, not $\Delta_5(Z)$. Why?

\textit{Closure.} The BKM root grading is on $\La^{2,1}_{II}$, while
the additive Gritsenko lift produces $\Delta_5$ on $\La^{2,1}_I$
in its primitive form, with $\La^{2,1}_{II}\subset\tfrac12\La^{2,1}_I$.
The doubling $Z\mapsto 2Z$ identifies the two lattice gradings; the
prefactor $64=[q^{1/2}r^{1/2}s^{1/2}]\Delta_5$ is the
theta-constant normalization and is untouched by doubling. See
Remark~\ref{rem:bkm-doubling-Z}.

\subsection*{Doubling factor as power of two}

\textit{Claim attacked.} The factor $64=2^6$ is associated with a
power of two from a half-Hilbert orbit.

\textit{Closure.} The factor $64$ is the value of the leading
Fourier coefficient $c_\Delta(1,1,1)$ of $\Delta_5$, computed in
the Maass formula for the product of ten even theta
constants~\cite[(1.3)--(1.5)]{GN}. It is not a half-Hilbert orbit
count. The orbit-count factor $2^{12}=4096$ arises in the type-II
Pfaffian wall normal form $(O2)$ of
Chapter~\ref{chap:dirac-igusa-pro-object} as the size of the half-Hilbert
orbit on the local sign atlas; it is unrelated to the denominator
prefactor $64$. See Remark~\ref{rem:bkm-doubling-power-of-two}.

\subsection*{Simple-root system on $\La^{2,1}_{II}$}

\textit{Claim attacked.} A different basis of three vectors of square
$2$ in $\La^{2,1}_{II}$ would yield a different BKM algebra.

\textit{Closure.} Nikulin's reflection-group
theorem~\cite{NIKULIN} pins
$W^{(2)}(\La^{2,1}_{II})$ as a normal subgroup of index $6$ in
$\Orth(\La^{2,1})_+$, and the Gritsenko--Nikulin chamber
$\Poly_{II}$ has fixed walls $\delta_1,\delta_2,\delta_3$
with Gram matrix~\eqref{eq:bkm-cartan-matrix}. Any other choice of
three primitive type-II square-$2$ vectors lies in the same Weyl
orbit as some permutation of $(\delta_1,\delta_2,\delta_3)$. The BKM
algebra depends on the equivalence class of the Cartan matrix, not
on the specific basis vectors.

\subsection*{Weyl vector $\rho$}

\textit{Claim attacked.} The Weyl vector might be
$\rho'=\delta_1+\delta_2+\delta_3$ rather than
$\tfrac12(\delta_1+\delta_2+\delta_3)$.

\textit{Closure.} The Weyl vector is fixed by the condition
$(\rho,\delta_i)=-(\delta_i,\delta_i)/2=-1$ for $i=1,2,3$. The
displayed value $\rho=\tfrac12(\delta_1+\delta_2+\delta_3)$
satisfies this and lies in $(\La^{2,1}_{II})^*\cap\Cone(\La^{2,1})_+$;
$\rho'=\delta_1+\delta_2+\delta_3$ would have $(\rho',\delta_i)=-2$,
violating the condition. The factor $\tfrac12$ is forced. See
Proposition~\ref{prop:bkm-weyl-vector}.

\subsection*{Real-root norm convention}

\textit{Claim attacked.} The convention $(\delta_i,\delta_i)=2$
might be exchanged with $(\delta_i,\delta_i)=1$ on a non-even
lattice.

\textit{Closure.} The lattice $\La^{2,1}=\HypPlan\oplus[2]$ is even
by~\eqref{eq:bkm-lambda21-form}. The type-II sublattice $\La^{2,1}_{II}$
is also even (a sublattice of an even lattice). Hence
$(\delta,\delta)\in 2\Z$ for every $\delta\in\La^{2,1}_{II}$. The
choice $(\delta,\delta)=2$ is forced for primitive square-norm-$2$
vectors. The convention $(\delta_i,\delta_i)=1$ does not occur.

\subsection*{Second Serre bracket compatibility (Gram cocycle)}

\textit{Claim attacked.} The Gram cocycle $B(c,c')$ on the upstairs
lattice $\Gamma_X^{\mathrm{phys}}$ might be compatible with the BKM
bracket only after normal ordering.

\textit{Closure.} The Gram cocycle $B(c,c')$ is the source-side
quadratic form on the upstairs lattice; its compatibility with the
BKM bracket on $\La^{2,1}_{II}$ is mediated by the normal-ordered
homomorphism $\overline\Pi_X\colon\widehat\Gamma_X\to
\Gamma_{\mathrm{gram}}$ of Chapter~\ref{ch:k3xe-setup},
which corrects the raw quadratic shadow
$\Pi_X(c+c')=\Pi_X(c)+\Pi_X(c')+B(c,c')$ to the additive map. The
present chapter, which is target-side, uses the Lorentzian Cartan
pairing on $\La^{2,1}_{II}$ directly; no Gram cocycle correction is
needed. The compatibility is precisely between
$\overline\Pi_X$ on the source and the BKM grading on the target,
as established in Chapter~\ref{ch:k3xe-setup}.

\subsection*{Alternative Cartan-matrix choice}

\textit{Claim attacked.} A symmetrized Cartan matrix
$\bigl(\begin{smallmatrix}1&-1&-1\\-1&1&-1\\-1&-1&1\end{smallmatrix}\bigr)$
would also satisfy the Borcherds--Kac axioms and might give the
same algebra.

\textit{Closure.} Borcherds--Kac--Moody algebras are classified up
to scalar normalization of the Cartan matrix by the equivalence class
of $A$. The matrix $A$ in~\eqref{eq:bkm-cartan-matrix} has integer
entries with $A_{ii}=2$, and is the symmetrized form fixed by the
even-lattice convention. A rescaled matrix
$\tfrac12 A$ with $A_{ii}=1$ would not be a generalized Cartan
matrix in the Kac--Moody sense (the condition $A_{ii}\in 2\Z_{>0}$
is part of the standard normalization in~\cite[Chapter 1]{KAC:1}).
The symmetric form $A$ is the unique normalization compatible with
the even lattice $\La^{2,1}_{II}$.

\subsection*{Cross-volume $\kappa_{\mathrm{BKM}}$ disagreement}

\textit{Claim attacked.} The Vol~I three-faces value $12$ might
disagree with this chapter's Borcherds-weight value $5$.

\textit{Closure.} The two values are the two indexings of the same
quantity:
$\kappa_{\mathrm{BKM}}^{(\mathrm{Borch})}=c_1(0)/2=5$ from
\cite[Theorem 13.3]{B} (Borcherds-weight indexing), and
$\kappa_{\mathrm{BKM}}^{(3\text{-faces})}=
2\kappa_{\mathrm{BKM}}^{(\mathrm{Borch})}+\chi(\mathcal O_X)+2=12$
on $X=K3\times E$ (three-faces indexing). The reconciliation
$12=2\cdot 5+0+2$ is a single identity, not a conflict. See
Remark~\ref{rem:bkm-cross-volume-reconciliation} and the Vol~III
universal-Borcherds-weight identity (Theorem~\ref{thm:bkm-kappa-universal}).

\subsection*{Determinant alone determines the algebra}

\textit{Claim attacked.} The denominator product
$64^{-1}\Delta_5(2Z)$ alone determines $\autcor$.

\textit{Closure.} False: a graded super vector space with the same
signed multiplicities and zero bracket gives the same denominator
product but is not a Lie superalgebra. The denominator identity
fixes only the signed Grothendieck class
(Proposition~\ref{prop:bkm-determinant-grothendieck}). The bracket,
the parity split, the radical, the simple primitives, and the
no-extra-relations theorem are required for a recognition theorem;
all of these are read on the BKM presentation of
Definition~\ref{def:bkm-relations}, not from the denominator product
alone. The recognition target $P_X\simeq\autcor$ requires the
compact-source data of Chapter~\ref{ch:pfaffian-dirac}.

\section{Coda: $\autcor$ as one of five views}
\label{sec:bkm-coda}

The denominator algebra $\autcor$ is one of five views of
$\Delta_5$: View~I (automorphic section), View~II (Borcherds--Kac--Moody
denominator, this chapter), View~III (compact protected Pfaffian,
conditional), View~IV (mirror discriminant, conjectural), View~V
(singular theta lift). The cross-level identities are:
\begin{itemize}
\item View~I $\leftrightarrow$ View~II: theorem,
\cite{Borcherds95,GNII}.
\item View~II $\leftrightarrow$ View~V: theorem,
\cite{B}.
\item View~I $\leftrightarrow$ View~V: theorem,
\cite{GN,B}.
\item View~II $\leftrightarrow$ View~III: \emph{conditional},
Chapter~\ref{ch:pfaffian-dirac}, under
$(D0),(O1),(O1)^+,(O2)$.
\item View~II $\leftrightarrow$ View~IV: \emph{conjectural},
Chapter~\ref{ch:view4-mirror}.
\end{itemize}

The closing identity of this chapter is the recorded denominator
identity
\begin{equation}
 \operatorname{den}(\autcor)\;=\;64^{-1}\Delta_5(2Z),
\end{equation}
the Borcherds-weight characteristic
\begin{equation}
 \kappa_{\mathrm{BKM}}(\Delta_5)\;=\;5,
\end{equation}
and the active multiplicity formula
\begin{equation}
 \smult(\alpha(n,l,m))\;=\;f(nm,l)\qquad((n,l,m)\in\Gamma_{\mathrm{act}}).
\end{equation}

These three identities are the load-bearing target-side data of
View~II. The chapter is closed.
